%==============================================================================
% Informe Proyecto Final - Estructuras de Datos y Analisis de Algoritmos
% UIS - Grupo B1/C1 - Mayo de 2026
% Simulador de Propagacion de Epidemia sobre Redes Sociales
%==============================================================================
\documentclass[spanish,notitlepage,letterpaper, 12pt]{article}

%---------------------------- Paquetes -----------------------------------------
\usepackage[utf8]{inputenc}
\usepackage[T1]{fontenc}
\usepackage[spanish,es-tabla]{babel}
\usepackage{amsmath}
\usepackage{amsfonts}
\usepackage{amssymb}
\usepackage{graphicx}
\usepackage{geometry}
\geometry{letterpaper}
\usepackage{epstopdf}
\usepackage{fancyhdr}
\usepackage{listings}
\usepackage{color}
\usepackage{xcolor}
\usepackage{booktabs}
\usepackage{array}
\usepackage{multirow}
\usepackage{longtable}
\usepackage{enumitem}
\usepackage{algorithm}
\usepackage{algpseudocode}
\usepackage{hyperref}
\usepackage{caption}
\usepackage{float}
\usepackage{tocloft}
\usepackage{tabularx}

%---------------------------- Ajustes de TOC -----------------------------------
\setlength{\cftsecnumwidth}{2.5em}
\setlength{\cftsubsecnumwidth}{3.2em}

%---------------------------- Colores y Listings -------------------------------
\definecolor{dkgreen}{rgb}{0,0.6,0}
\definecolor{gray}{rgb}{0.5,0.5,0.5}
\definecolor{mauve}{rgb}{0.58,0,0.82}
\lstset{frame=shadowbox,
  language=Java,
  aboveskip=3mm,
  belowskip=3mm,
  showstringspaces=false,
  columns=flexible,
  basicstyle={\small\ttfamily},
  numbers=left,
  numberstyle=\tiny\color{gray},
  keywordstyle=\color{blue},
  commentstyle=\color{dkgreen},
  stringstyle=\color{mauve},
  breaklines=true,
  breakatwhitespace=true,
  tabsize=3,
  rulesepcolor=\color{blue}
}

%---------------------------- Encabezado y Pie ---------------------------------
\pagestyle{fancy}
\lhead{}
\chead{\bfseries Proyecto Final EDA - Grupo B1/C1}
\rhead{Mayo de 2026}
\lfoot{\it Laura Galvis}
\cfoot{Universidad Industrial de Santander}
\rfoot{\thepage}
\voffset = -0.25in
\textheight = 8.0in
\textwidth = 6.5in
\oddsidemargin = 0.in
\headheight = 20pt
\headwidth = 6.5in
\renewcommand{\headrulewidth}{0.5pt}
\renewcommand{\footrulewidth}{0.5pt}

%---------------------------- Hyperref -----------------------------------------
\hypersetup{
  colorlinks=true,
  linkcolor=black,
  citecolor=black,
  urlcolor=blue
}

%---------------------------- Pseudocodigo en espanol --------------------------
\renewcommand{\algorithmicrequire}{\textbf{Entrada:}}
\renewcommand{\algorithmicensure}{\textbf{Salida:}}
\renewcommand{\algorithmicif}{\textbf{si}}
\renewcommand{\algorithmicthen}{\textbf{entonces}}
\renewcommand{\algorithmicelse}{\textbf{si no}}
\renewcommand{\algorithmicend}{\textbf{fin}}
\renewcommand{\algorithmicfor}{\textbf{para}}
\renewcommand{\algorithmicdo}{\textbf{hacer}}
\renewcommand{\algorithmicwhile}{\textbf{mientras}}
\renewcommand{\algorithmicreturn}{\textbf{retornar}}

\begin{document}

\title{\textbf{Estructuras de datos y an\'alisis de algoritmos \\ Proyecto Final - Grupo B1/C1} \\
Simulador de propagaci\'on de epidemia sobre redes sociales: comparaci\'on de seis estrategias de vacunaci\'on}

\author{
\textbf{Joshep Jared Lizarazo Montero} -- 2250150 \\
\textbf{Juan Esteban Barajas Mantilla} -- 2250183 \\
\textbf{Juan Pablo Rueda Angarita} -- 2250160
}
\date{Mayo de 2026}

\vspace{20pt}
\maketitle
\tableofcontents
\newpage

%==============================================================================
\section{Introducci\'on}
%==============================================================================

\subsection{Antecedentes y origen del proyecto}

Este trabajo evoluciona directamente de un proyecto previo que desarrollamos en la asignatura de \textbf{Matem\'aticas Discretas}, en el cual simulamos la propagaci\'on de una epidemia sobre distintos tipos de redes sociales usando \textit{Python}, \textit{NetworkX} y \textit{matplotlib}. En aquella versi\'on, los nodos del grafo eran entidades simples sin atributos propios y nos concentramos en comparar topolog\'ias te\'oricas (Erd\H{o}s--R\'enyi, Watts--Strogatz, Barab\'asi--Albert).

La versi\'on actual, que dise\~namos espec\'ificamente para la asignatura de \textbf{Estructuras de Datos y An\'alisis de Algoritmos}, sube la dificultad del proyecto en tres dimensiones:

\begin{enumerate}
    \item \textbf{Nodos con peso real:} cada persona del grafo tiene atributos demogr\'aficos (edad, estrato socioecon\'omico, ocupaci\'on) que influyen directamente en la probabilidad de contagio de sus aristas. El grafo deja de ser una abstracci\'on neutra y se convierte en un modelo demogr\'afico colombiano.
    \item \textbf{Implementaci\'on en Java con arquitectura por capas:} abandonamos el ecosistema de librer\'ias de alto nivel y construimos estructuras propias (listas de adyacencia sobre \texttt{HashMap}, colas de prioridad, conjuntos para sincron\'ia) sobre una arquitectura de cuatro capas inspirada en Domain-Driven Design.
    \item \textbf{Estrategia h\'ibrida propia:} incorporamos un algoritmo de vacunaci\'on original que combina la posici\'on estructural del nodo en la red (centralidad de intermediaci\'on y grado) con su perfil de riesgo individual (edad y estrato).
\end{enumerate}

\subsection{Motivaci\'on}

Las epidemias no se propagan de manera uniforme: su velocidad y alcance dependen directamente de c\'omo est\'an conectadas las personas entre s\'i. Una misma enfermedad con la misma probabilidad de contagio puede generar un brote devastador en una red densamente conectada o extinguirse r\'apidamente en una red fragmentada. Esta observaci\'on, validada experimentalmente durante la pandemia de COVID-19, nos motiv\'o a abordar el problema desde la computaci\'on: si contamos con recursos limitados de vacunaci\'on (solo el 20\% de la poblaci\'on), \textit{\textquestiondown a qui\'enes conviene vacunar primero?}

Las estrategias cl\'asicas (priorizar adultos mayores, vacunar al azar, vacunar a trabajadores de salud) responden a criterios m\'edicos o log\'isticos. Las estrategias \textit{estructurales} basadas en teor\'ia de grafos (priorizar \textit{hubs}, puentes inter-comunidad o focos densos) responden a criterios topol\'ogicos. En este proyecto las evaluamos todas dentro de un mismo banco de pruebas y proponemos una estrategia adicional que las fusiona.

\subsection{Objetivos}

\paragraph{Objetivo general.} Dise\~nar e implementar en Java un simulador de propagaci\'on epid\'emica sobre una red social generada con patrones demogr\'aficos colombianos reales y determinar computacionalmente cu\'al estrategia de vacunaci\'on minimiza el impacto del brote bajo restricciones de recursos.

\paragraph{Objetivos espec\'ificos.} Como equipo nos propusimos:
\begin{enumerate}[label=\arabic*.]
    \item Implementar un grafo dirigido y ponderado con listas de adyacencia usando \texttt{HashMap<Persona, List<Contacto>{}>}.
    \item Modelar la propagaci\'on mediante el modelo SIRV (Susceptible $\rightarrow$ Infectado $\rightarrow$ Recuperado $\rightarrow$ Vacunado).
    \item Generar autom\'aticamente una poblaci\'on con distribuci\'on demogr\'afica colombiana realista (edades, estratos socioecon\'omicos y ocupaciones) a partir de datos del DANE.
    \item Implementar seis estrategias de vacunaci\'on: aleatoria, por \textit{hubs}, por \textit{betweenness}, dentro de comunidades, Dijkstra Ponderado e h\'ibrida.
    \item Incorporar un mecanismo de actualizaci\'on din\'amica de pesos en tiempo de ejecuci\'on mediante eventos autom\'aticos activados por umbral de infectados.
    \item Visualizar la propagaci\'on en tiempo real usando GraphStream, con nodos coloreados seg\'un el estado SIRV.
    \item Comparar estad\'isticamente las seis estrategias y determinar la m\'as efectiva mediante un puntaje compuesto.
\end{enumerate}


\subsection{Alcance y restricciones}

Nuestro proyecto cumple las restricciones impuestas por la r\'ubrica del curso:

\begin{itemize}[leftmargin=*]
    \item \textbf{Equipo de tres integrantes.}
    \item \textbf{M\'as de un caso de prueba:} evaluamos la red principal de $N = 300$ nodos sobre tres modos de ejecuci\'on: individual por estrategia, comparativo de las seis estrategias simult\'aneamente, y experimento por lotes con 10\,000 grafos independientes por estrategia (60\,000 simulaciones totales).
    \item \textbf{Creaci\'on del grafo a trav\'es de archivo plano:} adem\'as de la generaci\'on autom\'atica, dise\~namos el sistema para aceptar la red desde archivos CSV (\texttt{personas\_red.csv} y \texttt{contactos\_red.csv}).
    \item \textbf{Mecanismo propio de actualizaci\'on de pesos:} dise\~namos un sistema de tres capas que recompone el peso efectivo de cada arista en cada turno: (i) eventos NPI por umbral (ALERTA\_LEVE, CUARENTENA, LOCKDOWN), (ii) reacci\'on local de vigilancia y (iii) fatiga social por declive prolongado de la epidemia.
    \item \textbf{Coherencia de la informaci\'on:} la red modela patrones reales colombianos extra\'idos del DANE; las restricciones m\'edicas (20\% de recursos, 15\% de pacientes cero) reflejan escenarios pl\'ausibles de salud p\'ublica.
    \item \textbf{Resultado claro:} cada simulaci\'on produce m\'etricas cuantitativas (pico de infectados, duraci\'on del brote, total de afectados, porcentaje de contenci\'on, $R_0$ estimado) y un veredicto rankeado de las estrategias seg\'un un score compuesto.
\end{itemize}

\subsection{Estructura del informe}

El cap\'itulo~\ref{sec:problema} formaliza el problema. El cap\'itulo~\ref{sec:marco} desarrolla el marco te\'orico (teor\'ia de grafos y modelo SIRV). El cap\'itulo~\ref{sec:informacion} describe la informaci\'on real que utilizamos. El cap\'itulo~\ref{sec:diseno} expone el dise\~no por capas, las estructuras de datos que elegimos y el proceso de generaci\'on de la red. El cap\'itulo~\ref{sec:estrategias} -- el m\'as extenso -- detalla las seis estrategias de vacunaci\'on, incluyendo nuestro aporte propio. El cap\'itulo~\ref{sec:eventos} cubre el mecanismo de actualizaci\'on de pesos exigido por la r\'ubrica. El cap\'itulo~\ref{sec:flujo} integra los componentes en el flujo end-to-end. El cap\'itulo~\ref{sec:analisis} presenta las m\'etricas y el score compuesto. El cap\'itulo~\ref{sec:resultados} muestra los resultados experimentales. El cap\'itulo~\ref{sec:originalidad} destaca nuestros aportes originales. Las conclusiones est\'an en el cap\'itulo~\ref{sec:conclusiones}. Los anexos contienen los enlaces a recursos (repositorio, video, diapositivas), el glosario de m\'etricas y las instrucciones de ejecuci\'on.

\newpage
%==============================================================================
\section{Descripci\'on del problema}
\label{sec:problema}
%==============================================================================

\subsection{Problema a resolver}

\begin{quote}
\itshape
Las epidemias se propagan de manera distinta seg\'un la estructura social de una comunidad y las caracter\'isticas individuales de sus integrantes. Buscamos modelar c\'omo una enfermedad se expande a trav\'es de una red social representada como un grafo, donde cada persona tiene atributos demogr\'aficos propios como edad, estrato socioecon\'omico y ocupaci\'on. A partir de esto, evaluamos qu\'e estrategia de vacunaci\'on --- con recursos limitados al 20\% de la poblaci\'on --- logra contener el brote de forma m\'as eficaz, comparando enfoques estructurales cl\'asicos contra una nueva estrategia h\'ibrida que combina la posici\'on del nodo en la red con su perfil de riesgo individual.
\end{quote}

\begin{figure}[H]
\centering
\includegraphics[width=4.5in]{Figures/Representacion_SIRV.jpeg}
\caption{Representaci\'on esquem\'atica del problema: una red social con personas (nodos) cuyos estados epidemiol\'ogicos (Susceptible, Infectado, Recuperado, Vacunado) evolucionan en el tiempo. El grosor de cada arista refleja la probabilidad de contagio.}
\label{fig:problema_general}
\end{figure}

\subsection{Formulaci\'on espec\'ifica}

A partir del enunciado anterior, el problema computacional que abordamos consiste en:

\begin{enumerate}
    \item Modelar la red social como un grafo dirigido y ponderado $G = (V, E, w)$, donde:
    \begin{itemize}
        \item $V$ es el conjunto de personas con sus atributos demogr\'aficos.
        \item $E \subseteq V \times V$ es el conjunto de aristas dirigidas que representan relaciones de contagio potencial.
        \item $w: E \rightarrow [0.05, 0.95]$ asigna a cada arista una probabilidad de contagio derivada de los atributos del nodo origen.
    \end{itemize}
    \item Simular la propagaci\'on usando el modelo SIRV en tiempo discreto (turnos).
    \item Implementar y comparar \textbf{seis estrategias} distintas de vacunaci\'on, cada una usando informaci\'on diferente del grafo (ninguna, grado, intermediaci\'on, densidad por comunidad, probabilidad de ruta, score compuesto).
    \item Producir un veredicto cuantitativo que rankee las estrategias seg\'un su efectividad para contener el brote.
\end{enumerate}

\subsection{Restricciones del modelo}

Para que la comparaci\'on entre estrategias sea v\'alida, impusimos las siguientes condiciones al sistema:

\begin{itemize}[leftmargin=*]
    \item \textbf{Recursos limitados:} solo podemos vacunar al \textbf{20\%} de la poblaci\'on. Es una restricci\'on dura del problema.
    \item \textbf{Pacientes cero al 15\%:} el brote arranca con un \textbf{15\% de infectados iniciales} elegidos antes de aplicar cualquier estrategia. Esto garantiza que las seis estrategias enfrenten exactamente la misma condici\'on inicial (los mismos pacientes cero, sobre la misma red, con la misma semilla aleatoria). Esta decisi\'on, que incorporamos en la versi\'on 8 del proyecto, es uno de nuestros aportes metodol\'ogicos y la discutimos en detalle en la secci\'on~\ref{sec:comparacion_justa}.
    \item \textbf{No reinfecci\'on:} una vez que un nodo pasa a Recuperado o Vacunado, no puede volver a Susceptible. Nuestro modelo asume inmunidad permanente.
    \item \textbf{M\'as de un caso de prueba:} evaluamos la red de $N = 300$ nodos en tres modos distintos (individual, comparativo y experimento por lotes con 10\,000 grafos por estrategia), tal como exige la r\'ubrica del curso.
    \item \textbf{Tiempo discreto:} la simulaci\'on procede por turnos. Dentro de cada turno aplicamos los cambios de estado de forma sincr\'onica (todos al final, no a medida que se procesan los nodos).
\end{itemize}

\subsection{Hip\'otesis de investigaci\'on}

Antes de implementar y ejecutar, planteamos como equipo siete hip\'otesis que el sistema deber\'ia poder contrastar. Estas hip\'otesis estructuran la fase experimental de nuestro proyecto.

\begin{table}[H]
\centering
\small
\caption{Hip\'otesis de investigaci\'on a verificar experimentalmente.}
\label{tab:hipotesis}
\begin{tabular}{c p{11.5cm}}
\toprule
\textbf{\#} & \textbf{Hip\'otesis} \\
\midrule
H1 & La vacunaci\'on por \textit{Hubs} es m\'as efectiva que la Aleatoria en redes con \textit{hubs} claros. \\
\addlinespace
H2 & La vacunaci\'on por \textit{Betweenness} es m\'as efectiva en redes con comunidades bien separadas. \\
\addlinespace
H3 & La estrategia H\'ibrida supera a \textit{Betweenness} puro en redes con disparidad de estratos socioecon\'omicos. \\
\addlinespace
H4 & En redes densas ($\sim 300$ nodos), los eventos por umbral tienen mayor impacto en la reducci\'on del pico. \\
\addlinespace
H5 & La vacunaci\'on Aleatoria puede ser comparable a \textit{Hubs} en redes homog\'eneas donde no hay \textit{hubs} dominantes. \\
\addlinespace
H6 & Comunidades supera a Aleatoria y a \textit{Hubs} en redes con \textit{clusters} intra-estrato densamente conectados. \\
\addlinespace
H7 & H\'ibrida supera a Comunidades en la red grande porque incorpora el perfil demogr\'afico individual, no solo la densidad del grupo. \\
\bottomrule
\end{tabular}
\end{table}

\newpage
%==============================================================================
\section{Marco te\'orico y modelo matem\'atico}
\label{sec:marco}
%==============================================================================

\subsection{Representaci\'on con teor\'ia de grafos}

\subsubsection{Grafo dirigido y ponderado}

Modelamos la red social como un grafo $G = (V, E, w)$ donde:

\begin{itemize}
    \item $V$ es el conjunto de v\'ertices (nodos), cada uno representando una persona.
    \item $E \subseteq V \times V$ es el conjunto de aristas dirigidas, donde $(u, v) \in E$ indica que la persona $u$ puede contagiar a la persona $v$.
    \item $w: E \rightarrow [0, 1]$ es la funci\'on de peso, que asigna a cada arista una probabilidad de contagio.
\end{itemize}

\subsubsection{Por qu\'e el grafo es dirigido}

Adoptamos la convenci\'on \texttt{agregarContacto(A, B, p)} con el significado ``A puede contagiar a B con probabilidad $p$''. El peso de la arista $A \rightarrow B$ se calcula usando los \textbf{factores de riesgo del nodo $A$} (el que contagia). Si queremos que $B$ tambi\'en pueda contagiar a $A$, agregamos expl\'icitamente la arista inversa $B \rightarrow A$ con los factores de $B$. Esto permite que, por ejemplo, un trabajador de salud (factor de ocupaci\'on 1.40) tenga aristas salientes m\'as pesadas que un jubilado (factor 0.80), aunque ambos tengan el mismo n\'umero de contactos.

\subsubsection{Atributos del nodo}

Cada nodo $v \in V$ tiene un vector de atributos demogr\'aficos:
\[
v = (\text{id}, \text{edad}, \text{estrato}, \text{ocupaci\'on}, \text{estado}_{\text{SIRV}})
\]

Estos atributos no son decorativos: influyen directamente en el peso de las aristas adyacentes al nodo.

\subsubsection{Funci\'on de peso de arista}

El peso de una arista $(u \rightarrow v)$ lo calculamos como:
\begin{equation}
w(u, v) = \text{clamp}\big( \, p_{\text{base}} \cdot f_{\text{edad}}(u) \cdot f_{\text{estrato}}(u) \cdot f_{\text{ocup}}(u), \; 0.05, \; 0.95 \, \big)
\label{eq:peso}
\end{equation}

donde:
\begin{align*}
p_{\text{base}} &\in [0.1, 0.4] \quad \text{(probabilidad base de contagio de la enfermedad)} \\
f_{\text{edad}}(u) &= 1.0 + \frac{\text{edad}(u)}{100} \quad \text{(personas mayores m\'as vulnerables)} \\
f_{\text{estrato}}(u) &= 1.0 + 0.1 \cdot (6 - \text{estrato}(u)) \quad \text{(estrato bajo $\Rightarrow$ mayor exposici\'on)}
\end{align*}

y $f_{\text{ocup}}(u)$ se especifica en la Tabla~\ref{tab:factor_ocupacion}.

\begin{table}[H]
\centering
\caption{Factor de ocupaci\'on $f_{\text{ocup}}(u)$.}
\label{tab:factor_ocupacion}
\begin{tabular}{lc}
\toprule
\textbf{Ocupaci\'on} & \textbf{Factor} \\
\midrule
Trabajador de salud & 1.40 \\
Trabajador informal & 1.30 \\
Estudiante & 1.20 \\
Empleado formal & 1.00 \\
Jubilado & 0.80 \\
\bottomrule
\end{tabular}
\end{table}

Normalizamos el resultado con la operaci\'on \textit{clamp} para que $w(u,v) \in [0.05, 0.95]$, evitando aristas con probabilidad 0 (nunca contagia) o 1 (siempre contagia), ambas situaciones poco realistas.

\subsection{Modelo epidemiol\'ogico SIRV}

El modelo SIRV (Susceptible -- Infectado -- Recuperado -- Vacunado) es una extensi\'on del cl\'asico modelo SIR de epidemiolog\'ia matem\'atica, adaptado a simulaci\'on discreta sobre grafos.

\subsubsection{Estados posibles}

Cada nodo en cada turno $t$ se encuentra en exactamente uno de cuatro estados:

\begin{table}[H]
\centering
\caption{Estados del modelo SIRV.}
\label{tab:estados_sirv}
\begin{tabular}{cl p{8.5cm}}
\toprule
\textbf{S\'imbolo} & \textbf{Estado} & \textbf{Descripci\'on} \\
\midrule
S & Susceptible & No ha sido infectado. Puede contraer la enfermedad. \\
I & Infectado & Tiene la enfermedad. Puede contagiar a sus vecinos S. \\
R & Recuperado & Super\'o la enfermedad. No puede reinfectarse ni contagiar. \\
V & Vacunado & Inmune desde el inicio. No puede infectarse. \\
\bottomrule
\end{tabular}
\end{table}

\subsubsection{Transiciones entre estados}

Las transiciones ocurren una vez por turno (tiempo discreto):
\begin{align*}
S \rightarrow I & \quad \text{con probabilidad } w(\text{infectado}, \text{susceptible}) \\
I \rightarrow R & \quad \text{cuando } \text{d\'iasInfectado} \geq \text{d\'iasRecuperaci\'on} \\
S \rightarrow V & \quad \text{antes de la simulaci\'on (estrategia de vacunaci\'on)}
\end{align*}

No existen las transiciones $I \rightarrow S$, $R \rightarrow S$ ni $V \rightarrow S$: el modelo asume inmunidad permanente.

\begin{figure}[H]
\centering
\includegraphics[width=4.5in]{Figures/Modelo_SIRV.jpg}
\caption{Diagrama de transiciones del modelo SIRV. Los estados R y V son absorbentes: no hay reinfecci\'on en este modelo.}
\label{fig:transiciones_sirv}
\end{figure}

\subsubsection{Ecuaci\'on de probabilidad de transici\'on por turno}

Para cada nodo $v$ con estado $S$ en el turno $t$, la probabilidad de pasar a $I$ en el turno $t+1$ es:
\begin{equation}
P\big( v \in I \text{ en } t+1 \big) = 1 - \prod_{\substack{u \in N(v) \\ \text{estado}(u,t) = I}} \big(1 - w(u, v)\big)
\label{eq:prob_transicion}
\end{equation}

donde $N(v)$ es el conjunto de vecinos de $v$ cuyos nodos apuntan hacia $v$ con aristas de contagio.

En la implementaci\'on discreta se simplifica: por cada vecino $u$ infectado, se genera un n\'umero aleatorio $r \in [0, 1)$. Si $r < w(u, v)$, el nodo $v$ se infecta en el siguiente turno.

\subsubsection{N\'umero reproductivo b\'asico $R_0$}

El $R_0$ estima cu\'antas personas nuevas infecta en promedio un nodo infectado durante su periodo de infecci\'on:
\begin{equation}
R_0 \approx \langle k \rangle \cdot p \cdot d
\end{equation}
donde $\langle k \rangle$ es el grado promedio de la red, $p$ es la probabilidad promedio de contagio por contacto y $d$ son los d\'ias promedio de infecci\'on. Si $R_0 > 1$ el brote se expande; si $R_0 < 1$ se extingue. Las estrategias de vacunaci\'on buscan reducir $R_0$ por debajo de 1.

\subsection{Tipos de redes sociales}

En el proyecto de Matem\'aticas Discretas se estudiaron cuatro topolog\'ias de red. En el proyecto actual la red se genera con el \textbf{modelo de comunidades} (\textit{Stochastic Block Model}), el m\'as realista para comunidades humanas colombianas. Se describen las cuatro como marco te\'orico:

\paragraph{Red Aleatoria (Erd\H{o}s--R\'enyi).} Cada par de nodos se conecta con probabilidad constante $p$, independientemente de cualquier estructura social. Genera redes homog\'eneas sin \textit{hubs} ni comunidades claras.

\paragraph{Red Small-World (Watts--Strogatz).} La mayor\'ia de nodos est\'an conectados localmente (como en una comunidad real), pero existe un peque\~no porcentaje $\beta$ de conexiones aleatorias de largo alcance. Esto genera el fen\'omeno de ``seis grados de separaci\'on'': la distancia promedio entre cualquier par de nodos es muy corta a pesar de que cada nodo tiene pocos vecinos.

\paragraph{Red Scale-Free (Barab\'asi--Albert).} Algunos nodos concentran una cantidad desproporcionada de conexiones (\textit{hubs}), mientras que la mayor\'ia tiene muy pocas. La distribuci\'on de grados sigue una ley de potencias: $P(k) \sim k^{-\gamma}$. Mecanismo de formaci\'on: nuevos nodos se conectan con mayor probabilidad a los que ya tienen m\'as conexiones (\textit{preferential attachment}).

\paragraph{Red de Comunidades (Stochastic Block Model).} \textbf{\textit{Topolog\'ia usada en este proyecto.}} Los nodos se agrupan en comunidades (familias, barrios). La probabilidad de conexi\'on dentro de una comunidad ($p_{\text{in}}$) es mucho mayor que entre comunidades ($p_{\text{out}}$), con $p_{\text{in}} \gg p_{\text{out}}$. Modela fielmente las redes sociales humanas.

\begin{figure}[H]
\centering
\includegraphics[width=6in]{Figures/ChatGPT Image 26 may 2026, 01_51_24 p.m..png}
\caption{Las cuatro topolog\'ias cl\'asicas de redes sociales: (a) Erd\H{o}s--R\'enyi (aleatoria), (b) Watts--Strogatz (small-world), (c) Barab\'asi--Albert (scale-free) y (d) Stochastic Block Model (comunidades). Esta \'ultima es la usada en el proyecto.}
\label{fig:tipos_redes}
\end{figure}

\subsection{Centralidades sobre grafos}

Dos m\'etricas de centralidad sustentan las estrategias estructurales:

\paragraph{Centralidad de grado.} Cantidad de aristas que salen (o entran a) un nodo:
\begin{equation}
\deg(v) = |N(v)|
\end{equation}
Los nodos con grado alto son \textbf{hubs}: si se infectan, contagian a muchos vecinos simult\'aneamente.

\paragraph{Centralidad de intermediaci\'on (\textit{betweenness}).} Fracci\'on de los caminos m\'as cortos entre pares de nodos que pasan por un nodo dado:
\begin{equation}
C_B(v) = \sum_{s \neq v \neq t} \frac{\sigma(s, t \mid v)}{\sigma(s, t)}
\label{eq:betweenness}
\end{equation}
donde $\sigma(s,t)$ es el n\'umero de caminos m\'as cortos entre $s$ y $t$, y $\sigma(s,t \mid v)$ es cu\'antos de esos caminos pasan por $v$. Un nodo con $C_B$ alto act\'ua de \textbf{puente} entre comunidades; si se vacuna, el virus queda atrapado en una comunidad.

\newpage
%==============================================================================
\section{Informaci\'on real utilizada}
\label{sec:informacion}
%==============================================================================

La r\'ubrica del proyecto exige que la informaci\'on utilizada para crear el grafo y las actualizaciones de pesos tenga \textbf{coherencia con el mundo real}. En esta secci\'on documentamos las fuentes y los criterios que usamos.

\subsection{Distribuciones demogr\'aficas (DANE)}

Generamos las personas con distribuciones basadas en datos del Departamento Administrativo Nacional de Estad\'istica (DANE) de Colombia.

\subsubsection{Distribuci\'on de edades}

Basada en la pir\'amide poblacional colombiana publicada por el DANE. Mayor concentraci\'on entre 15 y 35 a\~nos. A mayor edad, mayor vulnerabilidad biol\'ogica (mayor peso de contagio).

\begin{table}[H]
\centering
\caption{Distribuci\'on de edades implementada en \texttt{GeneradorPoblacion}.}
\label{tab:dist_edades}
\begin{tabular}{lc}
\toprule
\textbf{Rango de edad} & \textbf{Proporci\'on} \\
\midrule
0 -- 14 a\~nos (ni\~nos) & 24\% \\
15 -- 35 a\~nos (adultos j\'ovenes, m\'axima interacci\'on social) & 38\% \\
36 -- 59 a\~nos (adultos, interacci\'on laboral) & 26\% \\
60 -- 90 a\~nos (mayores, mayor vulnerabilidad) & 12\% \\
\bottomrule
\end{tabular}
\end{table}

\begin{figure}[H]
\centering
\includegraphics[width=4.5in]{Figures/images.populationpyramid.png}
\caption{Pir\'amide poblacional colombiana (DANE). Las edades centrales (15--35 a\~nos) concentran la mayor parte de la poblaci\'on y la mayor movilidad social.}
\label{fig:piramide}
\end{figure}

\subsubsection{Distribuci\'on de estratos}

Los estratos bajos en Colombia concentran a la mayor\'ia de la poblaci\'on. Adem\'as, generan m\'as aristas (hacinamiento, transporte p\'ublico masivo) con mayor peso de contagio.

\begin{table}[H]
\centering
\caption{Distribuci\'on de estratos socioecon\'omicos.}
\label{tab:dist_estratos}
\begin{tabular}{cc}
\toprule
\textbf{Estrato} & \textbf{Proporci\'on} \\
\midrule
1 & 30\% \\
2 & 30\% \\
3 & 30\% \\
4 & 5\% \\
5 & 3\% \\
6 & 2\% \\
\bottomrule
\end{tabular}
\end{table}

\subsubsection{Distribuci\'on de ocupaciones}

La ocupaci\'on se correlaciona con la edad:

\begin{itemize}
    \item Personas menores de 18 a\~nos $\rightarrow$ siempre \texttt{"estudiante"}.
    \item Personas mayores de 65 a\~nos $\rightarrow$ siempre \texttt{"jubilado"}.
    \item Resto $\rightarrow$ estudiante 35\%, informal 25\%, empleado 25\%, salud 10\%, jubilado 5\%.
\end{itemize}

\subsection{Factores de riesgo derivados}

A partir de los atributos demogr\'aficos, cada nodo expone tres factores num\'ericos que determinan el peso de sus aristas salientes:

\paragraph{Factor de edad.}
\[
f_{\text{edad}}(u) = 1.0 + \frac{\text{edad}(u)}{100}
\]
Mayor edad $\Rightarrow$ mayor vulnerabilidad biol\'ogica.

\paragraph{Factor de estrato.}
\[
f_{\text{estrato}}(u) = 1.0 + 0.1 \cdot (6 - \text{estrato}(u))
\]

\begin{table}[H]
\centering
\caption{Factor de estrato $f_{\text{estrato}}$.}
\label{tab:factor_estrato}
\begin{tabular}{cc l}
\toprule
\textbf{Estrato} & \textbf{Factor} & \textbf{Justificaci\'on} \\
\midrule
1 & 1.50 & Hacinamiento, transporte p\'ublico masivo \\
2 & 1.40 & Densidad de contacto alta \\
3 & 1.30 & Densidad media \\
4 & 1.20 & Contacto moderado \\
5 & 1.10 & Contacto bajo \\
6 & 1.00 & Menor exposici\'on al contacto masivo \\
\bottomrule
\end{tabular}
\end{table}

\paragraph{Factor de ocupaci\'on.} Ya presentado en la Tabla~\ref{tab:factor_ocupacion}.

\subsection{C\'alculo del peso de cada arista}

Antes de agregar cualquier arista, calculamos su \texttt{probContagio} aplicando la ecuaci\'on (\ref{eq:peso}). El siguiente listado en Java muestra nuestra implementaci\'on real:

\begin{lstlisting}[language=Java, caption={C\'alculo del peso de arista en \texttt{GeneradorPoblacion}.}]
// probBase configurable (en [0.1, 0.4])
public static double calcularProbContagio(Persona origen, double probBase) {
    double peso = probBase
            * origen.factorEdad()
            * origen.factorEstrato()
            * origen.factorOcupacion();
    return Math.max(0.05, Math.min(0.95, peso));
}
\end{lstlisting}

\paragraph{Ejemplo num\'erico.} Una enfermera de 50 a\~nos, estrato 2:
\begin{align*}
f_{\text{edad}} &= 1.0 + 50/100 = 1.50 \\
f_{\text{estrato}} &= 1.0 + 0.1 \cdot (6 - 2) = 1.40 \\
f_{\text{ocup}} &= 1.40 \quad \text{(salud)} \\
w &= \text{clamp}(0.20 \times 1.50 \times 1.40 \times 1.40, \; 0.05, \; 0.95) \\
  &= \text{clamp}(0.588, \; 0.05, \; 0.95) = 0.59
\end{align*}
Esa arista tiene casi un 60\% de probabilidad de contagio por turno, frente al 0.20 base de un empleado joven de estrato 6 sin ning\'un factor amplificador.

\subsection{Casos de prueba y carga desde archivo plano}

La r\'ubrica exige (i) m\'as de un caso de prueba y (ii) que el mecanismo de creaci\'on del grafo se realice a trav\'es de archivos planos. Nuestro proyecto cumple ambos requisitos:

\begin{itemize}[leftmargin=*]
    \item \textbf{Red principal ($N = 300$ nodos):} representa una comunidad urbana diversa con mezcla de estratos 1--6, ocupaciones variadas y m\'ultiples vecindarios conectados. Sobre esta misma red ejecutamos los tres modos de operaci\'on: individual por estrategia, comparativo de las seis simult\'aneamente, y experimento por lotes con 10\,000 grafos independientes por estrategia (60\,000 simulaciones totales).
\end{itemize}

Ambas modalidades (generaci\'on autom\'atica y carga desde CSV) est\'an disponibles. El formato es el siguiente:

\begin{lstlisting}[language={}, caption={Formato de \texttt{personas\_red.csv}.}]
id,nombre,edad,estrato,ocupacion
P001,Ana Torres,28,3,Docente
P002,Luis Herrera,35,2,Comerciante
P003,Maria Gomez,22,4,Estudiante
P004,Carlos Ruiz,45,3,Medico
...
\end{lstlisting}

\begin{lstlisting}[language={}, caption={Formato de \texttt{contactos\_red.csv}.}]
origen,destino,probContagio
P001,P002,0.55
P001,P005,0.50
P001,P011,0.45
P002,P003,0.60
...
\end{lstlisting}

La clase \texttt{CargadorRedCSV} acepta dos variantes de \texttt{personas} para retrocompatibilidad: el formato de 5 columnas (\texttt{id, nombre, edad, estrato, ocupacion}) y el formato moderno de 4 columnas (\texttt{id, edad, estrato, ocupacion}). En el formato moderno la columna \texttt{nombre} se ignora autom\'aticamente.

\newpage
%==============================================================================
\section{Dise\~no de la soluci\'on: estructuras de datos y arquitectura}
\label{sec:diseno}
%==============================================================================

\subsection{Arquitectura por capas}

Nuestro proyecto sigue una arquitectura de cuatro capas inspirada en Domain-Driven Design, adaptada a la complejidad del curso. Cada capa tiene una responsabilidad \'unica y solo depende de la capa inmediatamente inferior.

\begin{figure}[H]
\centering
\includegraphics[width=4.5in]{Figures/ChatGPT Image 26 may 2026, 01_54_13 p.m..png}
\caption{Imagen de la estructura del proyecto.}
\label{fig:arquitectura}
\end{figure}

\begin{table}[H]
\centering
\caption{Responsabilidades por capa.}
\label{tab:capas}
\small
\begin{tabular}{lp{10cm}}
\toprule
\textbf{Capa} & \textbf{Responsabilidad} \\
\midrule
\texttt{presentation} & Interacci\'on con el usuario: UI Swing (FlatDarkLaf), ventanas GraphStream, fallback por consola. \\
\texttt{application} & Orquesta los casos de uso: servicios, DTOs, \textit{commands}. \\
\texttt{domain} & N\'ucleo del problema: modelos, \textit{value objects} y algoritmos puros (independientes de framework). \\
\texttt{infrastructure} & Detalles t\'ecnicos: persistencia CSV/PDF, generaci\'on de poblaci\'on, c\'alculo de estad\'isticas. \\
\bottomrule
\end{tabular}
\end{table}

\subsection{Estructuras de datos elegidas}

\subsubsection{Lista de adyacencia sobre \texttt{HashMap}}

Representamos el grafo internamente con:

\begin{lstlisting}[language=Java, caption={Estructura interna de \texttt{RedSocial}.}]
public class RedSocial {
    private final Map<Persona, List<Contacto>> adyacencia = new HashMap<>();
    // ... metodos: agregarPersona, agregarContacto, getContactos,
    //              getTodasLasPersonas, getTodosLosContactos,
    //              getGrado, setearPacienteCero, ...
}
\end{lstlisting}

\paragraph{Justificaci\'on.} Elegimos lista de adyacencia (y no matriz) porque la red es \textbf{dispersa}: con $N = 300$ personas cada una tiene en promedio 5-15 contactos, por lo que la lista ocupa $O(N + M)$ mientras que una matriz ocupar\'ia $O(N^2) = 90\,000$ celdas, la mayor\'ia vac\'ias. Adem\'as, el \texttt{HashMap} permite acceder a los contactos de una persona en tiempo esperado $O(1)$.

\paragraph{Aspecto cr\'itico:} las clases \texttt{Persona} deben implementar \texttt{equals()} y \texttt{hashCode()} consistentes basados en el campo \texttt{id}. De lo contrario, dos instancias con el mismo \texttt{id} ser\'ian tratadas como personas distintas dentro del mapa.

\subsubsection{Cola de prioridad}

Usamos \texttt{PriorityQueue} en dos contextos:
\begin{itemize}
    \item \texttt{DijkstraPonderado}: como \textbf{max-heap} ordenado por probabilidad acumulada del camino. Permite extraer en cada iteraci\'on el nodo con mayor probabilidad alcanzada, implementando Dijkstra en su variante multiplicativa.
    \item \texttt{VacunacionHibrida}: como ranking de candidatos a vacunar, ordenados por \textit{score} compuesto descendente.
\end{itemize}

\subsubsection{Conjunto para sincron\'ia del SIRV}

En cada turno, los cambios de estado \textbf{no se aplican nodo a nodo} sino todos al final. Para esto usamos un \texttt{Set<Persona> nuevosInfectados} que acumula los candidatos antes de aplicar cualquier cambio. La raz\'on la discutimos en la secci\'on~\ref{sec:sincronia}.

\subsection{Modelo del dominio}

\subsubsection{\texttt{Persona} -- el nodo}

\begin{lstlisting}[language=Java, caption={Esqueleto de la clase \texttt{Persona} (paquete \texttt{domain.model}).}]
public class Persona {
    private final String id;        // unico - clave para HashMap
    private final int edad;
    private final int estrato;      // 1 a 6
    private final String ocupacion; // "salud","informal","estudiante","empleado","jubilado"
    private EstadoSIRV estado;      // SUSCEPTIBLE por defecto
    private int diasInfectado;      // contador para recuperacion

    public double factorEdad()      { return 1.0 + (edad / 100.0); }
    public double factorEstrato()   { return 1.0 + 0.1 * (6 - estrato); }
    public double factorOcupacion() {
        return switch (ocupacion) {
            case "salud"      -> 1.40;
            case "informal"   -> 1.30;
            case "estudiante" -> 1.20;
            case "jubilado"   -> 0.80;
            default           -> 1.00; // empleado u otros
        };
    }
    @Override public boolean equals(Object o) { /* por id */ }
    @Override public int hashCode()           { /* por id */ }
}
\end{lstlisting}

\subsubsection{\texttt{Contacto} -- la arista}

Representa una arista dirigida y ponderada. Almacena dos valores de peso: \texttt{probContagioBase} (inmutable, calculado una sola vez al construir la red) y \texttt{probContagio} (el valor efectivo del turno actual, reescrito cada turno por \texttt{AjustadorPesosAdaptativo}). Ambos est\'an clampados a $[0.05, 0.95]$.

\begin{lstlisting}[language=Java, caption={Clase \texttt{Contacto} (paquete \texttt{domain.model}).}]
public class Contacto {
    private final Persona origen, destino;
    private final double probContagioBase; // inmutable - base demografica
    private double probContagio;           // efectivo del turno actual

    public Contacto(Persona origen, Persona destino, double prob) {
        this.origen = origen;
        this.destino = destino;
        this.probContagioBase = clamp(prob, 0.05, 0.95);
        this.probContagio = this.probContagioBase;
    }
    public double getProbContagioBase() { return probContagioBase; }
    public void setProbContagio(double p) {
        this.probContagio = clamp(p, 0.05, 0.95);
    }
    private static double clamp(double v, double lo, double hi) {
        return Math.max(lo, Math.min(hi, v));
    }
}
\end{lstlisting}

\subsubsection{\texttt{EventoEpidemiologico} -- intervenci\'on no farmac\'eutica}

Representa un evento de salud p\'ublica activado cuando se cruza un umbral de infectados. La bandera \texttt{yaDisparado} garantiza que cada evento se dispara exactamente una vez. Cuando se activa, \texttt{GestorEventos} acumula su factor en \texttt{factorAcumuladoEventos}; el ajuste real de los pesos lo realiza \texttt{AjustadorPesosAdaptativo} cada turno a partir de ese factor. Se desarrolla en detalle en la secci\'on~\ref{sec:eventos}.


\subsection{Generaci\'on de la red: las seis fases del \texttt{GeneradorPoblacion}}

No construimos la red al azar puro. La hacemos pasar por \textbf{seis fases encadenadas} que imitan la estructura social real colombiana. El resultado es una red con propiedades de redes sociales reales: comunidades densas, \textit{hubs}, efecto small-world.

\begin{table}[H]
\centering
\small
\caption{Fases de construcci\'on de la red.}
\label{tab:fases}
\begin{tabular}{c p{4cm} p{8cm}}
\toprule
\textbf{Fase} & \textbf{Nombre} & \textbf{Descripci\'on} \\
\midrule
1 & Distribuciones demogr\'aficas (DANE) & Genera $N$ personas con edad, estrato y ocupaci\'on seg\'un las tablas DANE. \\
2 & \textit{Clusters} familiares & Grupos de 4-7 personas conectados como cliques con \texttt{probBase} 0.30-0.40. \\
3 & Vecindarios & 1-3 aristas moderadas (0.10-0.18) entre familias del mismo estrato. \\
4 & \textit{Hubs} comunitarios & 2-3 nodos de alta conectividad (mercado, iglesia) que conectan distintos \textit{clusters}. \\
4.5 & Conexiones de largo alcance & $N/15$ aristas aleatorias d\'ebiles (0.05-0.13) que producen el efecto small-world. \\
5 & Conectividad garantizada & BFS no dirigido para detectar componentes aisladas; se agregan aristas puente si las hay. \\
\bottomrule
\end{tabular}
\end{table}

\subsubsection{Fase 2 -- \textit{Clusters} familiares}

Las familias son \textbf{subgrafos casi completos} (cliques): todos los miembros est\'an conectados entre s\'i bidireccionalmente con \texttt{probBase} 0.30-0.40. Estas cliques son los ``n\'ucleos'' de contagio intra-familiar. Son exactamente las ``comunidades'' que detecta \texttt{VacunacionComunidades} midiendo densidad interna.

\subsubsection{Fase 3 -- Vecindarios}

Las familias de estratos similares se conectan entre s\'i con aristas moderadas. La probabilidad de conexi\'on decrece con la distancia de estrato:
\begin{itemize}
    \item distancia 1 (estratos adyacentes): 100\% de conexi\'on,
    \item distancia 2: 55\% de conexi\'on,
    \item distancia 3: 30\% de conexi\'on.
\end{itemize}
Estos puentes inter-familia son los que detecta \texttt{VacunacionBetweenness}.

\subsubsection{Fase 4 -- \textit{Hubs} comunitarios}

Creamos 2-3 nodos que simulan lugares de alta concentraci\'on (mercado, iglesia, transporte p\'ublico). Cada \textit{hub} conecta $N/10$ personas de \textbf{distintos clusters} con \texttt{probBase} moderada (0.08-0.15). Estos \textit{hubs} tienen grado mucho mayor que el promedio y son los nodos que \texttt{VacunacionHubs} y \texttt{VacunacionDijkstraPonderado} priorizan.

\subsubsection{Fase 4.5 -- Conexiones long-range (small-world)}

Agregamos $N/15$ aristas aleatorias entre nodos de cualquier parte de la red. Reducen el di\'ametro de la red, modelando ``conocidos lejanos'': compa\~neros de trabajo de otra ciudad, contactos digitales. Efecto clave: el virus puede llegar de cualquier comunidad a cualquier otra en muy pocos pasos.

\subsubsection{Fase 5 -- Conectividad garantizada}

Realizamos un BFS no dirigido sobre el grafo completo para detectar componentes conexas. Si quedan nodos no visitados (componentes hu\'erfanos), agregamos aristas bidireccionales puente para garantizar que la epidemia pueda alcanzar en teor\'ia a cualquier nodo desde cualquier otro.

\begin{figure}[H]
\centering
\includegraphics[width=6in]{Figures/Creacion_Grafo.jpeg}
\caption{Evoluci\'on del grafo a trav\'es de las seis fases de construcci\'on: desde $N$ nodos aislados (Fase 1) hasta una red social compleja con comunidades, \textit{hubs} y conexiones long-range (Fase 5).}
\label{fig:fases}
\end{figure}

\subsection{Propiedades emergentes de la red generada}

La red que produce nuestro \texttt{GeneradorPoblacion} tiene propiedades de redes sociales reales que explican por qu\'e las estrategias avanzadas superan a la aleatoria:

\begin{itemize}[leftmargin=*]
    \item \textbf{Estructura de comunidades:} los estratos forman grupos densos internamente y poco conectados entre s\'i. La aprovechan \texttt{VacunacionComunidades} y \texttt{VacunacionBetweenness}.
    \item \textbf{Hubs:} los nodos comunitarios (mercado, iglesia) tienen grado mucho mayor que el promedio. Los aprovechan \texttt{VacunacionHubs} y \texttt{VacunacionDijkstraPonderado}.
    \item \textbf{Efecto small-world:} el di\'ametro de la red es peque\~no gracias a las conexiones de largo alcance, aunque la mayor\'ia de contactos sean locales.
    \item \textbf{Pesos heterog\'eneos:} cada arista tiene su propia \texttt{probContagio} calibrada por el perfil del nodo origen. Los aprovechan \texttt{DijkstraPonderado} e \texttt{Hibrida}.
\end{itemize}

\newpage
%==============================================================================
\section{Estrategias de vacunaci\'on: los seis algoritmos}
\label{sec:estrategias}
%==============================================================================

Esta secci\'on, n\'ucleo del proyecto, presenta las seis estrategias de vacunaci\'on que implementamos. Todas vacunan exactamente el \textbf{20\% de los nodos susceptibles} antes de que inicie el brote; la diferencia entre ellas es \textbf{qu\'e nodos eligen}, y eso depende de \textit{qu\'e informaci\'on del grafo usan}.

\begin{table}[H]
\centering
\small
\caption{Resumen de las seis estrategias de vacunaci\'on.}
\label{tab:resumen_estrategias}
\begin{tabular}{c l l l p{5.5cm}}
\toprule
\textbf{\#} & \textbf{Estrategia} & \textbf{Nivel} & \textbf{Complejidad} & \textbf{Idea clave} \\
\midrule
1 & Aleatoria & Individuo & $O(N)$ & L\'inea base, sin informaci\'on de la red. \\
2 & Hubs & Individuo & $O(N \log N)$ & Vacuna los nodos de mayor grado. \\
3 & Betweenness & Individuo & $O(N(N+M))$ & Vacuna los puentes inter-comunidad (topol\'ogico). \\
4 & Comunidades & Grupo & $O(N+M)$ & Prioriza grupos de mayor densidad interna. \\
5 & Dijkstra Pond. & Individuo & $O(K N \log N)$ & Vacuna nodos en rutas de m\'axima prob. de contagio (Dijkstra multiplicativo). \\
6 & H\'ibrida $\star$ & Individuo & $O(N(N+M))$ & \textbf{Aporte propio}: combina estructura + perfil demogr\'afico. \\
\bottomrule
\end{tabular}
\end{table}



\subsection{Vacunaci\'on Aleatoria}

\paragraph{Base te\'orica.} Ninguna (l\'inea base).

\paragraph{Informaci\'on del grafo que usa.} Ninguna.

\paragraph{Complejidad.} $O(N)$.

\begin{algorithm}[H]
\caption{Vacunaci\'on Aleatoria}
\begin{algorithmic}[1]
\Require Red social $G$, semilla $s$
\State $L \gets$ lista de nodos en estado SUSCEPTIBLE
\State Mezclar $L$ con semilla $s$
\For{$i = 0$ \textbf{hasta} $\lfloor 0.2 \cdot |L| \rfloor - 1$}
    \State estado$(L[i]) \gets$ VACUNADO
\EndFor
\end{algorithmic}
\end{algorithm}

\paragraph{Para qu\'e sirve.} Es la \textbf{cota inferior}. Cualquier estrategia que use informaci\'on del grafo debe superarla para tener valor; si \textit{Betweenness} o \textit{Hubs} no la superan, hay un \textit{bug} o los par\'ametros est\'an mal.

\subsection{Vacunaci\'on por Hubs}

\paragraph{Base te\'orica.} Centralidad de grado.

\paragraph{Informaci\'on del grafo que usa.} Grado saliente de cada nodo.

\paragraph{Complejidad.} $O(N \log N)$ por el ordenamiento.

\begin{algorithm}[H]
\caption{Vacunaci\'on por Hubs}
\begin{algorithmic}[1]
\Require Red social $G$
\State $L \gets$ lista de nodos en estado SUSCEPTIBLE
\State Ordenar $L$ por $\deg(v)$ descendente
\For{$i = 0$ \textbf{hasta} $\lfloor 0.2 \cdot |L| \rfloor - 1$}
    \State estado$(L[i]) \gets$ VACUNADO
\EndFor
\end{algorithmic}
\end{algorithm}

\paragraph{Intuici\'on.} Vacunar a un \textit{hub} con grado 7 bloquea 7 aristas salientes de golpe; vacunar un nodo normal bloquea 1.

\paragraph{Fortaleza.} Muy efectiva en redes \textit{Scale-Free} donde los \textit{hubs} concentran la mayor\'ia de las conexiones.

\paragraph{Limitaci\'on.} No detecta nodos ``puente'' con grado moderado pero alta centralidad de intermediaci\'on. Un nodo con grado 3 que conecta dos comunidades de 50 personas puede ser m\'as estrat\'egico que un \textit{hub} con grado 10 dentro de una sola comunidad.

\subsection{Vacunaci\'on por Betweenness (algoritmo de Brandes)}

\paragraph{Base te\'orica.} Centralidad de intermediaci\'on, algoritmo de Brandes (2001).

\paragraph{Informaci\'on del grafo que usa.} Posici\'on global de cada nodo: todos los caminos m\'as cortos \textbf{topol\'ogicos} (por n\'umero de saltos).

\paragraph{Complejidad.} $O(N \cdot (N + M))$, dominante en redes dispersas.

\subsubsection{Qu\'e mide \textit{Betweenness}}

Repetimos la definici\'on de la ecuaci\'on (\ref{eq:betweenness}):
\[
C_B(v) = \sum_{s \neq v \neq t} \frac{\sigma(s,t \mid v)}{\sigma(s,t)}
\]
Un nodo con $C_B$ alto aparece en el camino m\'as corto entre muchos pares de nodos. En el contexto epidemiol\'ogico: \textbf{si se infecta, act\'ua de puente y lleva el virus de una comunidad a otra}.

\paragraph{Nota sobre ponderaci\'on.}
El algoritmo de Brandes implementado utiliza \textbf{BFS est\'andar no ponderado}: el ``camino m\'as corto'' se mide en n\'umero de saltos, \emph{no} en probabilidad acumulada de contagio. Esto significa que \textit{Betweenness} captura la posici\'on \textbf{estructural y topol\'ogica} del nodo (puentes entre comunidades) pero \emph{no} incorpora los pesos de las aristas. Si se quisiera integrar los pesos, la Fase~1 del algoritmo deber\'ia sustituirse por Dijkstra, elevando la complejidad a $O(N(M + N \log N))$. En nuestra arquitectura, la tarea de encontrar rutas de m\'axima probabilidad de contagio le corresponde a la estrategia \textit{Dijkstra Ponderado} (Secci\'on~\ref{sec:dijkstra}).

\subsubsection{Algoritmo de Brandes paso a paso}

Para cada nodo fuente $s$, se ejecutan dos fases.

\paragraph{Fase 1 -- BFS hacia adelante (cuenta caminos m\'inimos topol\'ogicos).}

\begin{algorithm}[H]
\caption{Fase 1 del algoritmo de Brandes desde la fuente $s$ (BFS no ponderado)}
\begin{algorithmic}[1]
\State $\sigma[s] \gets 1.0; \quad \sigma[v] \gets 0 \;\forall v \neq s$
\State $d[s] \gets 0; \quad d[v] \gets -1 \;\forall v \neq s$
\State $\text{pred}[v] \gets \emptyset \;\forall v$
\State $\text{pila} \gets \emptyset$
\State $Q \gets$ cola FIFO con $s$
\While{$Q$ no vac\'ia}
    \State $v \gets Q.\text{popleft}()$; apilar $v$ en $\text{pila}$
    \ForAll{vecino $w$ de $v$}
        \If{$d[w] = -1$}
            \State $d[w] \gets d[v] + 1$
            \State $Q.\text{push}(w)$
        \EndIf
        \If{$d[w] = d[v] + 1$}
            \State $\sigma[w] \gets \sigma[w] + \sigma[v]$
            \State $\text{pred}[w].\text{add}(v)$
        \EndIf
    \EndFor
\EndWhile
\end{algorithmic}
\end{algorithm}

\paragraph{Fase 2 -- Retropropagaci\'on (acumula la contribuci\'on a la centralidad).}

\begin{algorithm}[H]
\caption{Fase 2 del algoritmo de Brandes}
\begin{algorithmic}[1]
\State $\delta[v] \gets 0 \;\forall v$
\While{$\text{pila}$ no vac\'ia}
    \State $w \gets \text{pila}.pop()$
    \ForAll{$v \in \text{pred}[w]$}
        \State $\delta[v] \gets \delta[v] + \dfrac{\sigma[v]}{\sigma[w]} \cdot (1 + \delta[w])$
    \EndFor
    \If{$w \neq s$}
        \State $C_B[w] \gets C_B[w] + \delta[w]$
    \EndIf
\EndWhile
\end{algorithmic}
\end{algorithm}

La expresi\'on $\sigma[v]/\sigma[w]$ es la fracci\'on de caminos m\'inimos desde $s$ hasta $w$ que pasan por $v$. Multiplicado por $(1 + \delta[w])$ acumula tambi\'en la dependencia de todos los nodos alcanzados desde $w$. Brandes evita as\'i recalcular esto para cada par $(s,t)$, propagando la fracci\'on hacia atr\'as desde los nodos m\'as lejanos.

\paragraph{Fortaleza.} Muy efectiva en redes con comunidades para bloquear la propagaci\'on inter-comunidad.

\paragraph{Limitaci\'on.} Costosa computacionalmente en redes grandes. No considera los pesos de contagio (esa tarea corresponde a \textit{Dijkstra Ponderado}). En redes homog\'eneas su ventaja sobre \textit{Hubs} es peque\~na.

\subsection{Vacunaci\'on por Comunidades}

\paragraph{Base te\'orica.} Densidad interna de aristas por grupo.

\paragraph{Informaci\'on del grafo que usa.} Aristas internas de cada comunidad (estrato).

\paragraph{Complejidad.} $O(N + M)$ -- la m\'as eficiente de las estrategias estructurales.

\subsubsection{Estrato como proxy de comunidad}

En la red que produce nuestro \texttt{GeneradorPoblacion}, los \textit{clusters} familiares y vecinales se forman dentro del mismo estrato (las familias se conectan con otras familias de $|\text{estrato}_A - \text{estrato}_B| \leq 1$). Por eso usamos el estrato como un proxy estructuralmente v\'alido del grupo social colombiano.

\subsubsection{Definici\'on de densidad interna}

\begin{equation}
\rho(g) = \frac{\text{aristasInternas}(g)}{|g| \cdot (|g| - 1)}
\end{equation}
donde $\text{aristasInternas}(g)$ es el n\'umero de aristas cuyo origen Y destino pertenecen al mismo grupo $g$, y el denominador es el m\'aximo posible de aristas dirigidas sin lazos en $|g|$ nodos. Un grupo con $\rho = 0.8$ tiene el 80\% de todas las aristas posibles entre sus miembros: es un foco de contagio activo.

\begin{algorithm}[H]
\caption{Vacunaci\'on por Comunidades}
\begin{algorithmic}[1]
\Require Red social $G$
\State Agrupar nodos SUSCEPTIBLE por estrato: $\{g_1, g_2, \ldots, g_6\}$
\ForAll{$g \in \{g_1, \ldots, g_6\}$}
    \State Calcular $\rho(g)$
\EndFor
\State Ordenar grupos por $\rho$ descendente
\State $\text{cuota} \gets \lfloor 0.2 \cdot N_S \rfloor$ \Comment{$N_S$ = total de susceptibles}
\ForAll{$g$ en orden descendente de $\rho$}
    \State Ordenar miembros de $g$ por grado descendente
    \While{$\text{cuota} > 0$ y quedan miembros en $g$}
        \State Tomar el siguiente $v \in g$
        \State estado$(v) \gets$ VACUNADO
        \State $\text{cuota} \gets \text{cuota} - 1$
    \EndWhile
\EndFor
\end{algorithmic}
\end{algorithm}

\paragraph{Fortaleza.} Muy efectiva cuando las comunidades son densas y bien separadas. Corta focos intra-grupo. M\'as eficiente que \textit{Betweenness} e \textit{Hibrida}.

\paragraph{Limitaci\'on.} No bloquea los puentes inter-comunidad. Para eso \textit{Betweenness} es mejor. En redes homog\'eneas su resultado se acerca a Aleatoria.

\subsection{Vacunaci\'on por Dijkstra Ponderado}
\label{sec:dijkstra}

\paragraph{Base te\'orica.} Algoritmo de Dijkstra en su variante multiplicativa: maximiza el producto de probabilidades de contagio a lo largo de un camino.

\paragraph{Informaci\'on del grafo que usa.} Pesos reales de las aristas (probabilidades de contagio).

\paragraph{Complejidad.} $O(K \cdot N \log N)$ con $K = 5$ focos hipot\'eticos.

\subsubsection{El problema del camino de mayor contagio}

Dado el paciente cero, queremos encontrar el camino de \textbf{mayor probabilidad acumulada} de contagio hacia cualquier nodo del grafo, definido como:
\begin{equation}
\max_{p: u \rightsquigarrow v} \; P(u \rightarrow v) = \max_p \prod_{e_i \in p} w(e_i)
\end{equation}

Trabajamos en logaritmos (para evitar productos de flotantes muy peque\~nos):
\begin{equation}
\max_p \log P(u \rightarrow v) = \max_p \sum_{e_i \in p} \log w(e_i)
\end{equation}

Como $\log w(e_i) < 0$ para $w \in (0,1)$, maximizar la suma equivale a minimizar la suma de $-\log w(e_i)$, que es un problema cl\'asico de Dijkstra con pesos positivos. Por esto el algoritmo se denomina \textbf{Dijkstra Ponderado} (o Dijkstra multiplicativo): utiliza una cola de prioridad (\textit{max-heap} por probabilidad acumulada) que lo distingue fundamentalmente de BFS, el cual no considera los pesos de las aristas.

\subsubsection{Algoritmo}

\begin{algorithm}[H]
\caption{Dijkstra Ponderado -- camino de m\'axima probabilidad de contagio}
\begin{algorithmic}[1]
\Require Red social $G$, nodo origen $s$
\State $\text{prob}[s] \gets 1.0; \quad \text{prob}[v] \gets 0 \;\forall v \neq s$
\State $\text{pq} \gets$ max-heap con $(s, 1.0)$
\While{$\text{pq}$ no vac\'ia}
    \State $(u, \text{probU}) \gets \text{pq}.\text{pop}()$ \Comment{max por probabilidad acumulada}
    \If{$\text{probU} < \text{prob}[u]$} \textbf{continuar} \EndIf
    \ForAll{$(u \rightarrow v)$ en $G$}
        \State $\text{nuevaProb} \gets \text{probU} \cdot w(u, v)$
        \If{$\text{nuevaProb} > \text{prob}[v]$}
            \State $\text{prob}[v] \gets \text{nuevaProb}$
            \State $\text{pq}.\text{push}((v, \text{nuevaProb}))$
        \EndIf
    \EndFor
\EndWhile
\end{algorithmic}
\end{algorithm}

\subsubsection{Uso como estrategia de vacunaci\'on}

\texttt{VacunacionDijkstraPonderado} reutiliza este algoritmo:

\begin{enumerate}
    \item Selecciona $K = 5$ candidatos a paciente cero hipot\'etico: los nodos de mayor grado entre los susceptibles (los m\'as expuestos a ser focos iniciales).
    \item Desde cada candidato, ejecuta \texttt{DijkstraPonderado} hacia todos los dem\'as susceptibles, obteniendo el camino de mayor probabilidad acumulada.
    \item Para cada nodo intermedio del camino (excluye foco y destino), acumula como \textit{score} la probabilidad acumulada del camino completo. Los nodos que est\'an en muchas rutas cr\'iticas reciben \textit{score} alto.
    \item Vacuna el 20\% con mayor \textit{score} acumulado.
\end{enumerate}

\paragraph{Diferencia con \textit{Betweenness}.} Ambos identifican nodos estrat\'egicos en rutas de propagaci\'on, pero \textit{Betweenness} cuenta \textit{caminos m\'inimos topol\'ogicos} (saltos, sin importar el peso), mientras que \textit{Dijkstra Ponderado} usa la \textit{probabilidad real acumulada de la ruta} -- una m\'etrica m\'as alineada al riesgo epidemiol\'ogico real. Adem\'as, Dijkstra Ponderado eval\'ua solo $K = 5$ focos, lo que reduce el costo computacional respecto a $O(N(N+M))$ de \textit{Betweenness}.

\subsection{Vacunaci\'on H\'ibrida -- nuestro aporte propio $\star$}
\label{sec:hibrida}

\paragraph{Base te\'orica.} \textit{Score} compuesto: estructura de red + perfil individual.

\paragraph{Informaci\'on del grafo que usa.} Centralidad de intermediaci\'on, grado y atributos de cada nodo.

\paragraph{Complejidad.} $O(N(N+M))$, dominada por el c\'alculo de \textit{Betweenness}.

\subsubsection{Motivaci\'on}

Notamos que las estrategias anteriores consideran solo la posici\'on del nodo en la red. Sin embargo, en una comunidad real, la vulnerabilidad individual importa tanto como la conectividad: una persona mayor de estrato 1 con muchas conexiones deber\'ia tener prioridad sobre un joven de estrato 4 igualmente conectado, porque combina alto riesgo de propagar con alto riesgo de sufrir consecuencias graves. Por eso dise\~namos una estrategia que combina ambos criterios en una sola funci\'on objetivo.

\subsubsection{F\'ormula del \textit{score} compuesto}

Definimos el \textit{score} de cada nodo como una combinaci\'on lineal ponderada:
\begin{equation}
\boxed{
\text{score}(v) = \alpha \cdot C_B^{\text{norm}}(v) + \beta \cdot \text{riesgo\_edad}(v) + \gamma \cdot \text{vulnerabilidad\_estrato}(v) + \delta \cdot \deg^{\text{norm}}(v)
}
\label{eq:score_hibrido}
\end{equation}

con funciones auxiliares normalizadas y protegidas num\'ericamente:
\begin{align*}
C_B^{\text{norm}}(v) &= \dfrac{C_B(v)}{\max(C_B^{\max},\,\varepsilon)} \quad \in [0, 1] \\[6pt]
\deg^{\text{norm}}(v) &= \dfrac{\deg(v)}{\max(\deg^{\max},\,\varepsilon)} \quad \in [0, 1] \\[6pt]
\text{riesgo\_edad}(v) &= \dfrac{\text{edad}(v)}{90} \quad \in [0, 1] \\[6pt]
\text{vulnerabilidad\_estrato}(v) &= \dfrac{7 - \text{estrato}(v)}{6} \quad \in [0, 1]
\end{align*}

donde $\varepsilon = 10^{-9}$ es una constante de estabilidad num\'erica que evita divisi\'on por cero en grafos donde todos los nodos tienen intermediaci\'on nula ($C_B^{\max} = 0$) o grado cero ($\deg^{\max} = 0$). En Java, esta protecci\'on se implementa como:

\begin{lstlisting}[language=Java, caption={Normalizaci\'on con protecci\'on ante divisi\'on por cero.}]
final double EPS = 1e-9;
double cbNorm  = (maxBetweenness > EPS) ? betweenness.get(v) / maxBetweenness : 0.0;
double degNorm = (maxGrado > 0) ? (double) red.getGrado(v) / maxGrado : 0.0;
\end{lstlisting}

\subsubsection{Justificaci\'on de los pesos}

\begin{table}[H]
\centering
\caption{Ponderaci\'on de nuestro \textit{score} h\'ibrido.}
\label{tab:pesos_hibrida}
\begin{tabular}{l c c p{5cm}}
\toprule
\textbf{Componente} & \textbf{Par\'ametro} & \textbf{Valor} & \textbf{Justificaci\'on} \\
\midrule
Centralidad de intermediaci\'on & $\alpha$ & 0.35 & Es la se\~nal m\'as informativa: detecta puentes entre comunidades, mucho m\'as \'util que el grado puro. \\
Riesgo por edad & $\beta$ & 0.25 & Adultos mayores son los m\'as vulnerables ante la enfermedad. \\
Vulnerabilidad por estrato & $\gamma$ & 0.25 & Estratos bajos enfrentan condiciones de mayor exposici\'on. \\
Centralidad de grado (norm) & $\delta$ & 0.15 & Complementa $\alpha$ con la informaci\'on local. \\
\bottomrule
\end{tabular}
\end{table}

Los pesos suman $1.0$ y los dise\~namos para que:
\begin{itemize}
    \item Los \textbf{factores de red} ($\alpha + \delta = 0.50$) midan \emph{cu\'anto puede propagar} el virus el nodo.
    \item Los \textbf{factores personales} ($\beta + \gamma = 0.50$) midan \emph{cu\'anto sufrir\'ia} el individuo si se infectara.
\end{itemize}
Con la divisi\'on 50/50 le damos peso igual a ``cortar cadenas de contagio'' y ``proteger a los m\'as vulnerables''. Tomamos esta decisi\'on de dise\~no para reflejar la doble misi\'on de salud p\'ublica.

\subsubsection{Algoritmo}

\begin{algorithm}[H]
\caption{Vacunaci\'on H\'ibrida}
\small
\begin{algorithmic}[1]
\Require Red social $G$, par\'ametros $\alpha, \beta, \gamma, \delta$, constante $\varepsilon = 10^{-9}$
\State Calcular $C_B(v)$ y $\deg(v)$ para todos los nodos \Comment{algoritmo de Brandes, BFS no ponderado}
\State $C_B^{\max} \gets \max_v C_B(v)$; \quad $\deg^{\max} \gets \max_v \deg(v)$
\State $\text{pq} \gets$ max-heap (vac\'ia)
\ForAll{$v$ con estado SUSCEPTIBLE}
    \State $\hat{C}_B \gets C_B(v) \;/\; \max(C_B^{\max}, \varepsilon)$
    \State $\hat{d} \gets \deg(v) \;/\; \max(\deg^{\max}, \varepsilon)$
    \State $s \gets \alpha \cdot \hat{C}_B + \beta \cdot \dfrac{\text{edad}(v)}{90} + \gamma \cdot \dfrac{7 - \text{estrato}(v)}{6} + \delta \cdot \hat{d}$
    \State $\text{pq}.\text{push}((v, s))$
\EndFor
\For{$i = 1$ \textbf{hasta} $\lfloor 0.2 \cdot N_S \rfloor$}
    \State $(v, s) \gets \text{pq}.\text{pop}()$
    \State estado$(v) \gets$ VACUNADO
\EndFor
\end{algorithmic}
\end{algorithm}

\subsubsection{Ejemplo num\'erico comparativo}

Consideremos dos personas susceptibles en la misma red:

\paragraph{Persona A:} $C_B^{\text{norm}} = 0.8$, edad $= 30$, estrato $= 1$, $\deg^{\text{norm}} = 0.5$.
\begin{align*}
\text{score}(A) &= 0.35 \times 0.8 + 0.25 \times (30/90) + 0.25 \times ((7-1)/6) + 0.15 \times 0.5 \\
&= 0.28 + 0.083 + 0.25 + 0.075 = 0.688
\end{align*}

\paragraph{Persona B:} $C_B^{\text{norm}} = 0.2$, edad $= 75$, estrato $= 2$, $\deg^{\text{norm}} = 0.3$.
\begin{align*}
\text{score}(B) &= 0.35 \times 0.2 + 0.25 \times (75/90) + 0.25 \times ((7-2)/6) + 0.15 \times 0.3 \\
&= 0.07 + 0.208 + 0.208 + 0.045 = 0.531
\end{align*}

La estrategia H\'ibrida prioriza A porque su posici\'on en la red (\textit{betweenness} y grado altos) tiene m\'as impacto sist\'emico, aunque B sea mayor y m\'as vulnerable individualmente. Esta decisi\'on solo es posible gracias a la combinaci\'on de los cuatro factores.

\paragraph{Hip\'otesis del aporte.} Esta estrategia deber\'ia superar a \textit{Betweenness} puro en redes con comunidades de distintos estratos, porque prioriza nodos que son \emph{simult\'aneamente} puentes entre comunidades Y perfiles de alto riesgo (hip\'otesis H3 y H7).

\subsection{Comparaci\'on completa de las estrategias}

\begin{table}[H]
\centering
\scriptsize
\caption{Comparaci\'on completa de las seis estrategias de vacunaci\'on.}
\label{tab:comparacion_estrategias}
\begin{tabularx}{\textwidth}{l X X l X X}
\toprule
\textbf{Estrategia} & \textbf{Base te\'orica} & \textbf{Info del grafo} & \textbf{Complejidad} & \textbf{Fortaleza} & \textbf{Limitaci\'on} \\
\midrule
Aleatoria & Ninguna & Ninguna & $O(N)$ & Referencia neutral, rapid\'isima & Ignora la topolog\'ia \\
\addlinespace
Hubs & Grado (local) & Grado saliente de cada nodo & $O(N \log N)$ & Muy efectiva con \textit{hubs} pronunciados & No detecta puentes inter-comunidad \\
\addlinespace
Betweenness & Intermediaci\'on topol\'ogica & Caminos m\'inimos por saltos & $O(N^2)$ & Detecta puentes cr\'iticos estructurales & No usa pesos de contagio reales \\
\addlinespace
Comunidades & Densidad intra-grupo & Aristas internas por estrato & $O(N+M)$ & Corta contagio intra-grupo, muy r\'apida & No bloquea puentes entre comunidades \\
\addlinespace
Dijkstra Pond. & Dijkstra multiplicativo & Pesos reales desde $K$ focos & $O(K N \log N)$ & M\'as eficiente que \textit{Betw.}; usa probabilidades reales de contagio & Solo eval\'ua $K$ focos hipot\'eticos \\
\addlinespace
H\'ibrida $\star$ & Score compuesto & \textit{Betw.} topol\'ogico + grado + atributos & $O(N^2)$ & Combina red + perfil individual & Pesos heur\'isticos \\
\bottomrule
\end{tabularx}
\end{table}

\newpage
%==============================================================================
\section{Mecanismo de actualizaci\'on din\'amica de pesos}
\label{sec:eventos}
%==============================================================================

La r\'ubrica del proyecto exige expl\'icitamente que cada equipo dise\~ne ``\textit{un mecanismo propio para actualizar informaci\'on en medio de la ejecuci\'on, como pesos de aristas}''. Nuestra implementaci\'on, completada en la versi\'on 11 del proyecto, introduce \textbf{tres mecanismos complementarios} que se componen multiplicativamente cada turno para determinar el peso efectivo de cada arista:

\begin{equation}
\boxed{
w_{\text{ef}}(u{\to}v,\,t) = \mathrm{clamp}\!\Big(\, w_{\text{base}}(u{\to}v)\cdot f_{\text{eventos}}(t)\cdot \mathrm{vig}(v,t)\cdot \mathrm{fatiga}(t),\; 0.05,\; 0.95 \,\Big)
}
\label{eq:peso_efectivo}
\end{equation}

Los tres factores parten siempre del \texttt{probContagioBase} inmutable de la arista, lo que elimina el problema de acumulaci\'on de errores entre turnos. A continuaci\'on detallamos cada capa.

\subsection{Capa 1 -- Intervenciones No Farmac\'euticas, NPI (\texttt{GestorEventos})}

\subsubsection{Concepto}

En una epidemia real, los gobiernos reaccionan cuando la cantidad de infectados supera ciertos umbrales: primero recomiendan medidas b\'asicas, luego restringen la movilidad y finalmente confinan a la poblaci\'on. A estas medidas se les llama \textit{Intervenciones No Farmac\'euticas} (NPI). En nuestro c\'odigo modelamos el efecto conjunto de todos los NPI como un factor global acumulativo $f_{\text{eventos}}$ que disminuye cada vez que se cruza un umbral.

\subsubsection{Los tres eventos del sistema}

\begin{table}[H]
\centering
\caption{Eventos epidemiol\'ogicos autom\'aticos y su efecto sobre $f_{\text{eventos}}$.}
\label{tab:eventos}
\begin{tabular}{l c c c}
\toprule
\textbf{Evento} & \textbf{Umbral $\theta$} & \textbf{Factor $\lambda$} & \textbf{Reducci\'on efectiva} \\
\midrule
\texttt{ALERTA\_LEVE} & 30\% infectados & $\times 0.80$ & --20\% \\
\texttt{CUARENTENA} & 50\% infectados & $\times 0.50$ & --50\% \\
\texttt{LOCKDOWN} & 70\% infectados & $\times 0.20$ & --80\% \\
\bottomrule
\end{tabular}
\end{table}

Cada evento se dispara \textbf{exactamente una vez} (flag \texttt{yaDisparado}). El factor acumulado al dispararse los tres es:
\[
f_{\text{eventos}} = 0.80 \times 0.50 \times 0.20 = 0.08
\]

\subsubsection{Dise\~no del \texttt{GestorEventos} (v11)}

A diferencia de versiones anteriores en las que el \texttt{GestorEventos} iteraba sobre todas las aristas y llamaba directamente a \texttt{aplicarFactor()}, ahora \textbf{solo acumula el factor} en un campo interno y lo expone al ajustador:

\begin{algorithm}[H]
\caption{\texttt{GestorEventos.evaluar} -- acumula el factor NPI, ya no modifica aristas}
\begin{algorithmic}[1]
\Require Red social $G$, turno actual $t$, lista de eventos $\mathcal{E}$
\State porcentaje $\gets$ |nodos INFECTADO| / $|V|$
\State disparados $\gets$ lista vac\'ia
\ForAll{evento $e \in \mathcal{E}$ en orden ALERTA\_LEVE $\rightarrow$ CUARENTENA $\rightarrow$ LOCKDOWN}
    \If{$\neg e.\texttt{yaDisparado}$ \textbf{y} porcentaje $\geq e.\texttt{umbral}$}
        \State $f_{\text{eventos}} \gets f_{\text{eventos}} \times e.\texttt{factorMultiplicador}$
        \State $e.\texttt{disparar}()$ \Comment{marca \texttt{yaDisparado = true}}
        \State disparados.agregar($e$)
    \EndIf
\EndFor
\State \Return disparados \Comment{para que la UI muestre la alerta correspondiente}
\end{algorithmic}
\end{algorithm}

\subsubsection{Efecto acumulativo sobre $f_{\text{eventos}}$: ejemplo num\'erico}

\begin{table}[H]
\centering
\caption{Efecto acumulativo de los tres eventos sobre el factor NPI.}
\label{tab:efecto_acumulado}
\begin{tabular}{c l c c}
\toprule
\textbf{Turno} & \textbf{Evento disparado} & \textbf{C\'alculo} & \textbf{$f_{\text{eventos}}$} \\
\midrule
12 & \texttt{ALERTA\_LEVE} (30\%) & $1.00 \times 0.80$ & $0.80$ \\
18 & \texttt{CUARENTENA} (50\%) & $0.80 \times 0.50$ & $0.40$ \\
24 & \texttt{LOCKDOWN} (70\%) & $0.40 \times 0.20$ & $0.08$ \\
\bottomrule
\end{tabular}
\end{table}

El peso efectivo en el turno 24, con los dem\'as factores en 1.0 y \texttt{probBase} $= 0.50$, ser\'ia:
$w_{\text{ef}} = \mathrm{clamp}(0.50 \times 0.08, 0.05, 0.95) = 0.04$.

\subsection{Capa 2 -- Reacci\'on local (\texttt{AjustadorPesosAdaptativo})}
\label{sec:vigilancia}

\subsubsection{Concepto}

Cuando un nodo $v$ percibe que varios de sus vecinos directos est\'an infectados, aumenta su cautela: sale menos, usa protecci\'on, evita el contacto. Modelamos este fen\'omeno reduciendo la probabilidad de contagio en \emph{todas las aristas entrantes a} $v$ de forma proporcional a la fracci\'on de sus vecinos infectados:

\begin{equation}
\mathrm{vig}(v,t) = 1 - 0.60 \times \frac{|\{u : (u{\to}v)\in E \;\wedge\; \mathrm{estado}(u,t) = I\}|}{|\{u : (u{\to}v)\in E\}|}
\label{eq:vigilancia}
\end{equation}

\paragraph{Interpretaci\'on.}
\begin{itemize}
    \item Si el 100\% de los vecinos entrantes de $v$ est\'an infectados: $\mathrm{vig}(v) = 1 - 0.60 = 0.40$ (reducci\'on m\'axima del 60\%).
    \item Si ning\'un vecino est\'a infectado: $\mathrm{vig}(v) = 1.00$ (sin efecto).
\end{itemize}

\paragraph{Propiedades clave.}
\begin{itemize}[leftmargin=*]
    \item Se recalcula \textbf{desde cero} cada turno: no es acumulativa ni persistente.
    \item Es \textbf{completamente local}: $v$ reacciona solo a lo que observa en su entorno inmediato, independientemente del estado global de la epidemia.
    \item Afecta las aristas \emph{entrantes} a $v$: la cautela de $v$ reduce la probabilidad de que cualquier vecino $u$ infectado lo contagie.
\end{itemize}

\subsection{Capa 3 -- Fatiga social (\texttt{AjustadorPesosAdaptativo})}
\label{sec:fatiga}

\subsubsection{Concepto}

Cuando la epidemia lleva varios turnos sin crecer, la poblaci\'on tiende a relajar las precauciones: la gente sale m\'as, baja la guardia, incumple las restricciones. Modelamos este fen\'omeno como un aumento gradual de los pesos cuando $I(t)$ lleva $\geq 5$ turnos consecutivos sin crecer:

\begin{equation}
\mathrm{fatiga}(t) = \begin{cases}
1.0 & \text{si } \Delta < 0 \\[6pt]
1.0 + \min\!\left(0.30,\;\dfrac{\Delta}{10} \times 0.30\right) & \text{si } \Delta \geq 0
\end{cases}
\label{eq:fatiga}
\end{equation}

donde $\Delta = \mathrm{turnosDecreciendo}(t) - 5$. Cuando $I(t)$ vuelve a crecer, $\mathrm{turnosDecreciendo}$ se reinicia a 0 y $\mathrm{fatiga}(t) = 1.0$ de inmediato.

\paragraph{Ejemplo num\'erico.}
\begin{itemize}
    \item 5 turnos bajando: $\Delta = 0 \Rightarrow \mathrm{fatiga} = 1.00$ (umbral reci\'en alcanzado).
    \item 10 turnos bajando: $\Delta = 5 \Rightarrow \mathrm{fatiga} = 1.15$ (+15\%).
    \item 15+ turnos bajando: $\Delta \geq 10 \Rightarrow \mathrm{fatiga} = 1.30$ (+30\%, m\'aximo).
\end{itemize}

\paragraph{Interacci\'on con los dem\'as factores.} La fatiga ``empuja hacia arriba'' contrarrestando parcialmente las reducciones de los eventos NPI y de la vigilancia local, sin superar el tope de \textit{clamp} (0.95). Modela el \textbf{efecto memoria}: aunque el confinamiento a\'un tenga efecto ($f_{\text{eventos}} < 1$), las personas se relajan conforme el peligro parece menor, lo que puede provocar un repunte tard\'io de la epidemia.

\subsection{Integraci\'on: \texttt{AjustadorPesosAdaptativo}}
\label{sec:ajustador}

\subsubsection{Flujo de actualizaci\'on por turno}

\texttt{AjustadorPesosAdaptativo} es invocado por \texttt{SimulacionService} una vez por turno, despu\'es de que \texttt{GestorEventos} ha evaluado los umbrales:

\begin{algorithm}[H]
\caption{\texttt{AjustadorPesosAdaptativo.ajustar} -- recompone todos los pesos desde la base}
\begin{algorithmic}[1]
\Require Red $G$, factor $f_{\text{eventos}}$, historial $H$ de conteos SIRV
\State Construir $\mathrm{vig}[v]$ para todo $v \in V$ \Comment{ecuaci\'on~(\ref{eq:vigilancia})}
\State Calcular $\mathrm{fatiga}$ a partir de $H$ \Comment{ecuaci\'on~(\ref{eq:fatiga})}
\ForAll{arista $(u{\to}v) \in G$}
    \State $c.\texttt{setProbContagio}\!\big(\,c.\texttt{probContagioBase} \times f_{\text{eventos}} \times \mathrm{vig}[v] \times \mathrm{fatiga}\,\big)$
\EndFor
\State \Return \texttt{CambioFatiga} \Comment{\texttt{NINGUNO}, \texttt{INICIADA} o \texttt{REINICIADA}}
\end{algorithmic}
\end{algorithm}

\texttt{Contacto} almacena \texttt{probContagioBase} (campo \texttt{final}, calculado una sola vez al construir la red) y \texttt{probContagio} (valor efectivo del turno). El ajustador siempre parte de la base, por lo que no se acumulan errores entre turnos.

\subsubsection{Ejemplo num\'erico compuesto}

Arista P001$\to$P043 en el turno 22:
\begin{align*}
w_{\text{base}} &= 0.55 \quad \text{(trabajador de salud, estrato 2, 45 a\~nos)} \\
f_{\text{eventos}} &= 0.40 \quad \text{(ALERTA\_LEVE $\times$ CUARENTENA ya disparados)} \\
\mathrm{vig}(\text{P043}) &= 0.52 \quad \text{(60\% de vecinos infectados: $1 - 0.6 \times 0.8 = 0.52$)} \\
\mathrm{fatiga} &= 1.00 \quad \text{(la epidemia sigue creciendo)} \\
w_{\text{ef}} &= \mathrm{clamp}(0.55 \times 0.40 \times 0.52 \times 1.00,\; 0.05,\; 0.95) = 0.114
\end{align*}

Si en el turno 35 la epidemia lleva 8 turnos en declive y todos los NPI est\'an activos:
\begin{align*}
f_{\text{eventos}} &= 0.08, \quad \mathrm{vig} = 1.00, \quad \mathrm{fatiga} = 1.09 \\
w_{\text{ef}} &= \mathrm{clamp}(0.55 \times 0.08 \times 1.00 \times 1.09,\; 0.05,\; 0.95) = \mathrm{clamp}(0.048, 0.05, 0.95) = 0.05
\end{align*}

\subsubsection{Alertas en pantalla}

El m\'etodo retorna el enum \texttt{CambioFatiga \{NINGUNO, INICIADA, REINICIADA\}} para desacoplar el dominio de la presentaci\'on. \texttt{SimulacionService} usa ese valor junto con la lista de eventos disparados por \texttt{GestorEventos} para llamar a los m\'etodos de alerta de \texttt{GraficoSimulacion}:

\begin{itemize}[leftmargin=*]
    \item \textbf{Evento NPI:} banner naranja (\texttt{ALERTA\_LEVE}), naranja oscuro (\texttt{CUARENTENA}) o rojo (\texttt{LOCKDOWN}), visible 4~s.
    \item \textbf{Fatiga iniciada:} banner verde (``la poblaci\'on se relaja'').
    \item \textbf{Fatiga reiniciada:} banner dorado (``precauci\'on restaurada por rebote'').
\end{itemize}

\subsubsection{Historial de pesos en \texttt{VentanaResultados}}

Para que el usuario pueda observar c\'omo los tres factores afectan el peso de cada nodo, la pestaña \textbf{``Historial de pesos''} de \texttt{VentanaResultados} (solo en modos individual y comparativo) muestra una tabla nodo$\times$turno donde cada celda es el promedio de \texttt{probContagio} de las aristas salientes del nodo en ese turno. En el modo comparativo hay una sub-pesta\~na por cada una de las seis estrategias, lo que permite comparar visualmente c\'omo los distintos conjuntos de vacunados afectan la din\'amica de los pesos a lo largo del brote.



\subsection{Comparaci\'on justa: paciente cero antes de vacunar (v8)}
\label{sec:comparacion_justa}

Una de nuestras contribuciones metodol\'ogicas (incorporada en la versi\'on 8 del proyecto) es garantizar que las seis estrategias enfrenten exactamente la misma condici\'on inicial.

\paragraph{Problema previo.} En nuestra primera implementaci\'on vacun\'abamos PRIMERO y luego eleg\'iamos el paciente cero sobre los susceptibles restantes. Como cada estrategia vacuna nodos distintos, el conjunto de susceptibles cambiaba y, aun con la misma semilla, los infectados iniciales terminaban siendo distintos por estrategia. La comparaci\'on no era justa.

\paragraph{Cambio 1a -- orden.} Ahora ejecutamos \texttt{setearPacienteCero(\ldots)} \textbf{antes} de \texttt{vacunacionCommand.ejecutar(\ldots)}. El paciente cero se elige sobre la poblaci\'on completa e id\'entica (misma red + misma semilla); las seis estrategias arrancan con exactamente los mismos infectados. Como todos los algoritmos de vacunaci\'on filtran por \texttt{SUSCEPTIBLE}, los infectados quedan excluidos autom\'aticamente (nunca se vacuna a un nodo ya infectado).

\paragraph{Cambio 1b -- 15\% de infectados iniciales.} Aplicamos la constante \texttt{FRACCION\_INFECTADOS\_INICIALES = 0.15} en el \textit{setter} de \texttt{ConfiguracionDto}: con $N = 300$ generamos 45 infectados iniciales. La cuota de vacunaci\'on (20\%) la calculamos sobre los susceptibles \textbf{restantes} tras fijar el paciente cero. Este dise\~no garantiza condiciones iniciales id\'enticas para las seis estrategias.

\newpage
%==============================================================================
\section{Flujo de trabajo y unificaci\'on del sistema}
\label{sec:flujo}
%==============================================================================

\subsection{Flujo end-to-end}


El flujo completo de ejecuci\'on es:

\begin{enumerate}[leftmargin=*]
    \item \texttt{Main.java} arranca y decide entre Swing o consola.
    \item \texttt{VentanaMenuPrincipal} (o \texttt{ConsolaMenu}) recoge el modo y la configuraci\'on del usuario.
    \item \texttt{VentanaConfiguracion} construye un \texttt{ConfiguracionDto}.
    \item \texttt{IniciarSimulacionCommand} orquesta:
    \begin{enumerate}
        \item \texttt{GeneradorPoblacion.generar(N, semilla)} o \texttt{CargadorRedCSV.cargar(\ldots)} produce la \texttt{RedSocial}.
        \item \texttt{red.setearPacienteCero(15\% de N, random)} infecta los mismos nodos en las seis estrategias.
        \item \texttt{VacunacionService.aplicar(estrategia)} vacuna el 20\% de los SUSCEPTIBLES restantes.
        \item \texttt{SimulacionService.ejecutar(\ldots)} corre el bucle de turnos.
    \end{enumerate}
    \item Por cada turno:
    \begin{enumerate}
        \item \texttt{ModeloSIRV.simularTurno(red)}: recolecta nuevos infectados y nuevos recuperados; aplica todos los cambios sincr\'onicamente al final del turno.
        \item \texttt{GestorEventos.evaluar(red, turno)}: revisa umbrales; si se cruza uno, acumula el factor en \texttt{factorAcumuladoEventos} (ya no modifica aristas directamente).
        \item Registra el conteo $\{S, I, R, V\}$ del turno en el historial.
        \item \texttt{AjustadorPesosAdaptativo.ajustar(red, factorEventos, historial)}: recompone el peso efectivo de cada arista como $w_{\text{base}} \times f_{\text{eventos}} \times \mathrm{vig}(v) \times \mathrm{fatiga}$. Retorna \texttt{CambioFatiga} para la UI.
        \item \texttt{GraficoSimulacion.actualizarTurno(\ldots)}: actualiza colores en pantalla, la barra de turno SIRV y los banners de alerta si corresponde.
    \end{enumerate}
    \item Al final, \texttt{CalculadorEstadisticas} calcula las m\'etricas y \texttt{AnalisisComparativo} rankea las estrategias.
    \item \texttt{VentanaResultados} muestra: (a) tabla de m\'etricas con fila ganadora resaltada, (b) ranking con score compuesto, (c) curva $I(t)$ y (d) pesta\~na \textbf{``Historial de pesos''} (solo individual/comparativo) con la evoluci\'on del peso efectivo nodo$\times$turno.
    \item \texttt{ExportadorResultados} y \texttt{GeneradorReportePDF} producen los entregables.
\end{enumerate}

\subsection{Tres modos de operaci\'on}

\begin{table}[H]
\centering
\caption{Modos de simulaci\'on disponibles en la interfaz.}
\label{tab:modos}
\small
\begin{tabular}{l p{10cm}}
\toprule
\textbf{Modo} & \textbf{Descripci\'on} \\
\midrule
\textsc{Individual} & Corre una sola estrategia (elegida por el usuario) con visualizaci\'on GraphStream en una ventana. \\
\textsc{Comparativo} & Corre las seis estrategias sobre la \textbf{misma red} (mismos infectados iniciales) y las muestra en pesta\~nas separadas (\texttt{JTabbedPane}). \\
\textsc{Experimento por lotes} & Corre \textbf{10\,000 grafos distintos por estrategia} (6 $\times$ 10\,000 = 60\,000 grafos independientes), sin animaci\'on, y promedia los resultados. Veredicto estad\'isticamente robusto. \\
\textsc{Construcci\'on visual} & Anima la generaci\'on de la red fase por fase y arista por arista, con un color distinto por fase: clusters familiares, vecindarios, hubs, long-range y aristas puente. Implementado en \texttt{VisualizadorConstruccionRed} sobre \texttt{GeneradorPoblacion.generarConFases}. \\
\bottomrule
\end{tabular}
\end{table}

\subsection{Sincron\'ia del modelo SIRV}
\label{sec:sincronia}

En cada turno, los cambios \textbf{no se aplican nodo a nodo} sino todos al final. Esto es cr\'itico para la correcci\'on del modelo:

\paragraph{Sin sincron\'ia (incorrecto).} En el turno $t$, al iterar en orden \texttt{P001, P002, P003\ldots}, si \texttt{P001} (INFECTADO) contagia a \texttt{P002}, \texttt{P002} pasar\'ia a INFECTADO en el mismo turno $t$. Luego \texttt{P002}, ahora infectada, podr\'ia contagiar a \texttt{P003} dentro del mismo turno. El virus se propagar\'ia de forma irreal, dependiendo del orden de iteraci\'on.

\paragraph{Con sincron\'ia (correcto).} En el turno $t$, recolectamos todos los nodos que se van a infectar en un \texttt{Set<Persona> nuevosInfectados}, sin modificar el estado. Al finalizar la recolecci\'on, aplicamos todos los cambios simult\'aneamente. As\'i \texttt{P003} solo puede infectarse en el turno $t+1$, nunca en $t$.

\subsection{Visualizaci\'on en tiempo real}

Nuestro \texttt{GraficoSimulacion} envuelve la API de \texttt{GraphStream 2.0} para representar el grafo visualmente:
\begin{itemize}
    \item Crea un nodo visual por cada \texttt{Persona} y una arista visual por cada \texttt{Contacto}.
    \item Actualiza el color de cada nodo en cada turno seg\'un su \texttt{EstadoSIRV}: \textcolor{blue}{azul} para Susceptible, \textcolor{red}{rojo} para Infectado, \textcolor{green!50!black}{verde} para Recuperado, \textcolor{orange}{naranja} para Vacunado.
    \item El grosor de las aristas refleja \texttt{probContagio} efectiva del turno, que cambia cada turno con los tres factores adaptativos.
    \item Muestra una \textbf{barra de turno} debajo del grafo con el conteo SIRV en vivo: \texttt{Turno N \quad S: .. \quad I: .. \quad R: .. \quad V: ..}.
    \item Muestra \textbf{banners de alerta} sobre el grafo cuando se dispara un evento NPI (naranja / naranja oscuro / rojo seg\'un el evento) o cuando cambia el estado de la fatiga social (verde al iniciar la relajaci\'on, dorado al restaurar la precauci\'on). El banner se oculta autom\'aticamente tras 4~s mediante un \texttt{javax.swing.Timer}.
\end{itemize}

En modo comparativo, una \texttt{VentanaComparativaTabs} aloja las seis simulaciones en pesta\~nas separadas, todas ejecutadas secuencialmente sobre la misma red.

\newpage
%==============================================================================
\section{An\'alisis comparativo y m\'etricas}
\label{sec:analisis}
%==============================================================================

\subsection{M\'etricas de comparaci\'on}

Al finalizar cada simulaci\'on, \texttt{CalculadorEstadisticas} calcula seis m\'etricas:

\begin{table}[H]
\centering
\small
\caption{M\'etricas calculadas al final de cada simulaci\'on.}
\label{tab:metricas}
\begin{tabular}{l p{6.5cm} p{3.5cm}}
\toprule
\textbf{M\'etrica} & \textbf{Definici\'on} & \textbf{Direcci\'on favorable} \\
\midrule
Pico m\'aximo de infectados $I_{\max}$ & $\max\{I(t) : t = 0, 1, \ldots, T\}$ & menor = mejor \\
Turno del pico $t_{\text{pico}}$ & $\arg\max_t I(t)$ & mayor = mejor (m\'as tarde) \\
Duraci\'on del brote $T_{\text{brote}}$ & $\min\{t : I(t) = 0\}$ & menor = mejor \\
Total de afectados & $R(T_{\text{brote}})$ & menor = mejor \\
Contenci\'on (\%) & $(N - \text{total\_afectados}) / N \times 100$ & mayor = mejor \\
$R_0$ estimado & promedio de nuevos infectados por nodo infectado en $t = 1, 2, 3$ & menor = mejor \\
\bottomrule
\end{tabular}
\end{table}

\subsection{Score compuesto -- determinaci\'on cuantitativa del ganador}

Nuestro \texttt{AnalisisComparativo} rankea las estrategias con un \textit{score} compuesto sobre las cinco m\'etricas principales (excluyendo $t_{\text{pico}}$ que es derivada):

\begin{equation}
\boxed{
\begin{aligned}
\text{score} = \; & 0.30 \cdot \text{norm\_inv}(I_{\max}) \\
            + \; & 0.15 \cdot \text{norm\_inv}(T_{\text{brote}}) \\
            + \; & 0.25 \cdot \text{norm\_inv}(\text{total\_afectados}) \\
            + \; & 0.20 \cdot \text{norm}(\text{contenci\'on\%}) \\
            + \; & 0.10 \cdot \text{norm\_inv}(R_0)
\end{aligned}
}
\label{eq:score_compuesto}
\end{equation}

donde $\text{norm\_inv}$ invierte las m\'etricas ``menor = mejor'' (pico, duraci\'on, afectados, $R_0$) para que en todos los casos un valor m\'as alto del t\'ermino signifique ``mejor''. La suma de los pesos es 1.0, por lo que $\text{score} \in [0, 1]$.

\paragraph{Justificaci\'on de los pesos.} Elegimos los pesos as\'i:

\begin{itemize}
    \item $0.30$ pico: presi\'on sobre el sistema de salud -- la m\'etrica de impacto m\'as visible.
    \item $0.25$ afectados: cu\'antos pasaron por la enfermedad -- define el costo humano total.
    \item $0.20$ contenci\'on: cu\'anta poblaci\'on nunca se infect\'o (incluye los vacunados).
    \item $0.15$ duraci\'on: brote corto = menos disrupci\'on social.
    \item $0.10$ $R_0$: velocidad inicial -- informativa, pero correlacionada con pico.
\end{itemize}

La salida es una \texttt{List<ScoreEstrategia>} ordenada de mayor a menor \textit{score}. El primer elemento es la estrategia ganadora con su desglose por componente y una justificaci\'on narrativa que generamos autom\'aticamente.

\subsection{Modo por lotes: promedio + victorias acumuladas}

En nuestro experimento por lotes, evaluamos cada estrategia sobre 10\,000 grafos distintos e independientes. Cada par (estrategia $s$, corrida $i$) usa una semilla \'unica $\text{base} + s \cdot 10000 + i$, produciendo 60\,000 grafos diferentes en total. Al terminar reportamos dos lecturas complementarias:

\paragraph{Lectura promedio.} Promediamos las m\'etricas de las 10\,000 corridas de cada estrategia (pico, duraci\'on, afectados, contenci\'on, $R_0$) y las rankeamos con el mismo \textit{score} compuesto del comparativo.

\paragraph{Lectura acumulada (victorias).} En cada corrida rankeamos las seis estrategias por \textit{score} y contamos una ``victoria'' para la mejor. Es una estimaci\'on Monte Carlo de qu\'e tan seguido cada estrategia resulta la mejor sobre grafos aleatorios. Con 10\,000 corridas por estrategia, el intervalo de confianza del porcentaje de victorias es $\pm 1\%$ (al 95\%), lo que hace el veredicto estad\'isticamente s\'olido.

\subsection{Glosario de m\'etricas en los reportes}

Nuestros reportes PDF (tanto el individual/comparativo como el de lotes) incluyen al final una p\'agina de \textbf{Glosario de m\'etricas -- Gu\'ia de interpretaci\'on}: una tabla que explica cada variable, qu\'e mide y qu\'e valores se consideran favorables desde el punto de vista de salud p\'ublica. Reproducimos el glosario completo en el Anexo~\ref{anexo:glosario} de este informe.

\newpage
\newpage
%==============================================================================
\section{Resultados experimentales}
\label{sec:resultados}
%==============================================================================

\subsection{Configuración experimental}

Realizamos las pruebas sobre la red principal de $N = 300$ nodos en sus tres modos
de operación:

\begin{table}[H]
\centering
\caption{Configuración experimental.}
\label{tab:config_exp}
\small
\begin{tabular}{l p{8cm}}
\toprule
\textbf{Caso} & \textbf{Parámetros} \\
\midrule
Red principal (ciudad, $N = 300$) & $\text{probBase} = 0.20$,
  díasRecuperación $= 7$, turnos máximos $= 50$,
  15\% pacientes cero (45 nodos), 20\% vacunados. \\
Lote (estadístico) & 6 estrategias $\times$ 10\,000 grafos
  = 60\,000 simulaciones independientes, $N = 300$.
  Generado el 2026-05-26. \\
\bottomrule
\end{tabular}
\end{table}

\subsection{Resultados sobre la Red principal ($N = 300$ nodos)}

\begin{figure}[H]
\centering
\includegraphics[width=4.5in]{Figures/WhatsApp Image 2026-05-26 at 1.30.57 PM.jpeg}
\caption{Ventana principal del simulador (\texttt{VentanaMenuPrincipal}) con tema
FlatDarkLaf y selección de modo.}
\label{fig:ventana_menu}
\end{figure}

\begin{figure}[H]
\centering
\includegraphics[width=5in]{Figures/WhatsApp Image 2026-05-26 at 5.49.07 PM.jpeg}
\caption{Visualización en tiempo real con GraphStream: nodos coloreados por estado
SIRV, aristas con grosor proporcional a \texttt{probContagio}, y barra de turno
con conteo en vivo bajo el grafo.}
\label{fig:graphstream}
\end{figure}

\begin{figure}[H]
\centering
\includegraphics[width=5.5in]{Figures/WhatsApp Image 2026-05-26 at 5.49.28 PM.jpeg}
\caption{\texttt{VentanaResultados}: tabla de métricas por estrategia (con fila
ganadora resaltada en verde) y panel de ranking con el \textit{score} compuesto
y la justificación narrativa.}
\label{fig:ventana_resultados}
\end{figure}

\begin{figure}[H]
\centering
\includegraphics[width=5.5in]{Figures/WhatsApp Image 2026-05-26 at 5.50.31 PM.jpeg}
\caption{Página del PDF generado: curva $I(t)$ comparativa de las seis estrategias
superpuestas en un solo gráfico.}
\label{fig:pdf_curva}
\end{figure}

\begin{figure}[H]
\centering
\includegraphics[width=5.5in]{Figures/WhatsApp Image 2026-05-26 at 5.50.24 PM.jpeg}
\caption{Página del PDF: barras por métrica (pico, duración, afectados,
contención\,\%, $R_0$) con etiquetas numéricas y paleta por estrategia.}
\label{fig:pdf_barras}
\end{figure}

\begin{figure}[H]
\centering
\includegraphics[width=5.5in]{Figures/WhatsApp Image 2026-05-26 at 5.50.58 PM.jpeg}
\caption{Página del PDF: desglose apilado del \textit{score} compuesto. Cada
barra muestra qué componente contribuyó más al \textit{score} final de cada
estrategia.}
\label{fig:pdf_score}
\end{figure}

\subsection{Experimento por lotes — veredicto estadístico}
\label{sec:resultados_lote}

En nuestro experimento por lotes evaluamos las seis estrategias sobre 10\,000
grafos distintos cada una (60\,000 simulaciones en total). La
Tabla~\ref{tab:resultados_lote} resume los promedios de todas las métricas y las
victorias acumuladas.

\begin{table}[H]
\centering
\caption{Resultados del experimento por lotes ($N=300$, 10\,000 grafos por
estrategia, 60\,000 simulaciones totales). Métricas = promedio de las 10\,000
corridas. Victorias = corridas en que la estrategia obtuvo el mayor \textit{score}
compuesto. Fila destacada = estrategia ganadora.}
\label{tab:resultados_lote}
\small
\begin{tabular}{l c c c c c c c}
\toprule
\textbf{Estrategia} & \textbf{Pico} & \textbf{t-Pico} & \textbf{Dur.}
  & \textbf{Afectados} & \textbf{Cont.\%} & \textbf{$R_0$}
  & \textbf{Victorias} \\
\midrule
\textbf{HUBS}            & \textbf{131} & 6 & \textbf{26}
  & \textbf{180} & \textbf{40,0\%} & \textbf{2,06}
  & \textbf{6\,980 / 10\,000} \\
BETWEENNESS              & 150 & 6 & 27 & 203 & 32,3\% & 2,12
  & 1\,405 / 10\,000 \\
Dijkstra Ponderado       & 156 & 6 & 27 & 208 & 30,7\% & 2,15
  &    760 / 10\,000 \\
HÍBRIDA                  & 159 & 6 & 27 & 212 & 29,3\% & 2,15
  &    547 / 10\,000 \\
COMUNIDADES              & 169 & 6 & 26 & 216 & 28,0\% & 2,16
  &    291 / 10\,000 \\
ALEATORIA                & 187 & 6 & 24 & 227 & 24,3\% & 2,23
  &     17 / 10\,000 \\
\bottomrule
\multicolumn{8}{l}{\textbf{Ganadora:} HUBS — \textit{score} compuesto promedio
  0,355\,/\,1,000 — 6\,980 de 10\,000 corridas ganadas (69,8\%).}
\end{tabular}
\end{table}

\begin{figure}[H]
\centering
\includegraphics[width=5.5in]{Figures/WhatsApp Image 2026-05-26 at 5.50.31 PM.jpeg}
\caption{PDF del modo por lotes: victorias acumuladas por estrategia en las
10\,000 corridas independientes. HUBS domina con 6\,980 victorias (69,8\%),
seguida de BETWEENNESS con 1\,405 (14,1\%).}
\label{fig:pdf_lotes}
\end{figure}

La estrategia \textbf{HUBS} resultó ganadora con un margen amplio y consistente.
Su fortaleza principal proviene de la métrica de \textbf{contención}: al vacunar
los nodos de mayor grado, bloquea de golpe las aristas de mayor alcance en la red
de comunidades, impidiendo que el virus salte rápidamente de un \textit{cluster}
a otro. Con solo el 20\,\% de la población vacunada, logró que el 40\,\% de la
red nunca se infectara, frente al 24,3\,\% de la estrategia Aleatoria.

\subsection{Discusión y validación de hipótesis}

\begin{table}[H]
\centering
\footnotesize
\caption{Validación de hipótesis según el experimento por lotes (10\,000 corridas,
$N = 300$).}
\label{tab:validacion}
\begin{tabular}{c p{5.8cm} p{5.5cm}}
\toprule
\textbf{H\#} & \textbf{Hipótesis} & \textbf{Resultado observado} \\
\midrule
H1 & \textit{Hubs} supera a Aleatoria en redes con \textit{hubs} claros.
   & \textbf{Confirmada ampliamente.} HUBS: 40\,\% de contención vs.\ 24,3\,\%
     de Aleatoria; 6\,980 vs.\ 17 victorias. \\
\addlinespace
H2 & \textit{Betweenness} es más efectiva en redes con comunidades separadas.
   & \textbf{Confirmada parcialmente.} Betweenness queda $2.^{\circ}$ (1\,405
     victorias), superada por HUBS. \\
\addlinespace
H3 & Híbrida supera a \textit{Betweenness} con disparidad de estratos.
   & \textbf{No confirmada.} Híbrida queda $4.^{\circ}$ (547 victorias) detrás
     de Betweenness (1\,405) y Dijkstra Ponderado (760). \\
\addlinespace
H4 & Eventos por umbral tienen mayor impacto en redes densas ($N = 300$).
   & \textbf{Confirmada.} Los eventos CUARENTENA y LOCKDOWN se dispararon
     sistemáticamente antes que en redes pequeñas. \\
\addlinespace
H5 & Aleatoria es comparable a \textit{Hubs} en redes homogéneas.
   & \textbf{Refutada en esta topología.} La red de comunidades genera
     \textit{hubs} pronunciados que hacen decisiva la diferencia (6\,980 vs.\
     17 victorias). \\
\addlinespace
H6 & Comunidades supera a Aleatoria y \textit{Hubs} con \textit{clusters} densos.
   & \textbf{Parcialmente confirmada.} Supera a Aleatoria (291 vs.\ 17
     victorias), pero queda $5.^{\circ}$ en el ranking general. \\
\addlinespace
H7 & Híbrida supera a Comunidades en redes grandes.
   & \textbf{Confirmada.} Híbrida ($4.^{\circ}$, 547 victorias) supera a
     Comunidades ($5.^{\circ}$, 291 victorias), aunque ambas quedan lejos del
     podio en esta topología. \\
\bottomrule
\end{tabular}
\end{table}

\paragraph{Análisis del resultado inesperado en H3.}
La hipótesis H3 planteaba que la estrategia Híbrida superaría a Betweenness gracias
a su componente demográfica. El experimento por lotes muestra lo contrario: en la
red de $N = 300$ con modelo de comunidades, la posición estructural pura del nodo
(intermediación topológica) es más informativa que la combinación con el perfil
individual. Una explicación plausible es que en esta topología los nodos ``puente''
entre comunidades concentran también los perfiles de mayor riesgo (estratos bajos,
trabajadores informales), lo que hace que la señal demográfica de Híbrida no aporte
información adicional relevante sobre Betweenness. Este resultado abre la línea de
trabajo futuro de calibración automática de los pesos $\alpha, \beta, \gamma,
\delta$ en función de la topología de la red.

\paragraph{Dominancia de HUBS.}
El resultado más llamativo es la dominancia de HUBS con 6\,980 victorias sobre
10\,000 corridas (69,8\,\%). En el modelo de comunidades que genera nuestro
\texttt{GeneradorPoblacion}, los nodos ``mercado'' e ``iglesia'' (Fase 4 de la
construcción) concentran conexiones hacia múltiples \textit{clusters} de distintos
estratos. Vacunarlos equivale a cerrar los únicos canales de contagio inter-comunidad
de forma simultánea, lo que tiene un efecto desproporcionadamente grande sobre el
pico y el total de afectados. Este resultado valida de forma cuantitativa la hipótesis
H1 y es consistente con la literatura sobre redes \textit{scale-free}~\cite{barabasi}.
\newpage
%==============================================================================
\section{Originalidad del aporte}
\label{sec:originalidad}
%==============================================================================

La r\'ubrica del curso indica expl\'icitamente que se ``\textit{evaluar\'a la originalidad de la soluci\'on propuesta en comparaci\'on con los dem\'as grupos}''. En esta secci\'on resumimos nuestros cuatro aportes originales.

\subsection{Estrategia H\'ibrida por score compuesto}

La quinta estrategia (secci\'on~\ref{sec:hibrida}) es nuestra principal contribuci\'on. Sus elementos originales son:
\begin{itemize}[leftmargin=*]
    \item \textbf{Fusi\'on estructura+demograf\'ia} en una sola funci\'on objetivo, en lugar de elegir un solo criterio (grado, intermediaci\'on, densidad de comunidad).
    \item \textbf{Calibraci\'on de pesos 50/50} (red vs perfil individual) que justificamos con criterio de salud p\'ublica, no por optimizaci\'on \textit{post-hoc}.
    \item \textbf{Cuatro factores normalizados con protecci\'on num\'erica} ($\varepsilon$-guard) para evitar divisiones por cero y dominancias arbitrarias entre magnitudes.
\end{itemize}

\subsection{Modo por lotes con doble veredicto (promedio + victorias)}

Originalmente nuestro sistema solo soportaba el modo comparativo de una sola corrida, lo que daba veredictos inestables (sensibles a la semilla). Como equipo dise\~namos el modo por lotes (versi\'on 8) que aporta:
\begin{itemize}[leftmargin=*]
    \item Evaluaci\'on sobre 10\,000 grafos \textbf{distintos e independientes} por estrategia (60\,000 simulaciones totales), con semillas \'unicas por par (estrategia, corrida).
    \item Dos lecturas complementarias: \textbf{promedio} (consistencia) y \textbf{victorias acumuladas} (estimaci\'on Monte Carlo de cu\'al estrategia es la mejor con mayor frecuencia).
    \item Reporte PDF dedicado (\texttt{GeneradorReporteLotePDF}) con tabla extendida que incluye la columna ``Victorias''.
\end{itemize}

\subsection{Comparaci\'on justa: paciente cero antes de vacunar (v8)}

Identificamos que nuestra implementaci\'on inicial introduc\'ia un sesgo en la comparaci\'on (cada estrategia arrancaba con un conjunto de pacientes cero distinto, debido a que los eleg\'iamos \emph{despu\'es} de vacunar). Como equipo corregimos el orden:
\begin{enumerate}
    \item Generar la red.
    \item \textbf{Fijar el 15\% de pacientes cero} sobre la poblaci\'on completa con la misma semilla.
    \item Aplicar la estrategia de vacunaci\'on sobre los susceptibles restantes.
    \item Ejecutar la simulaci\'on.
\end{enumerate}
Esto garantiza que las seis estrategias enfrenten \textbf{exactamente la misma condici\'on inicial}, lo que hace que las comparaciones sean estad\'isticamente v\'alidas.

\subsection{Pesos din\'amicos adaptativos (v11)}

M\'as all\'a del requerimiento m\'inimo de ``un mecanismo de actualizaci\'on de pesos'', dise\~namos un sistema de \textbf{tres capas independientes} que operan simult\'aneamente en cada turno (ecuaci\'on~(\ref{eq:peso_efectivo})):

\begin{itemize}[leftmargin=*]
    \item \textbf{Reacci\'on local} (\texttt{AjustadorPesosAdaptativo}): cada nodo reduce la probabilidad de sus aristas entrantes seg\'un la fracci\'on de sus vecinos infectados. Es completamente local y se reinicia cada turno.
    \item \textbf{Fatiga social} (\texttt{AjustadorPesosAdaptativo}): cuando la epidemia lleva $\geq 5$ turnos sin crecer, los pesos aumentan gradualmente hasta +30\% modelando el relajamiento de precauciones; se reinicia cuando la epidemia rebota.
    \item \textbf{Retroalimentaci\'on visual}: banners de color en \texttt{GraficoSimulacion} hacen visible al usuario el momento exacto en que el modelo cambia el comportamiento de la red, tanto para los NPI como para los cambios de fatiga.
    \item \textbf{Historial de pesos}: la pesta\~na ``Historial de pesos'' de \texttt{VentanaResultados} (solo individual/comparativo) muestra una tabla nodo$\times$turno con la evoluci\'on del peso efectivo, permitiendo analizar c\'omo cada estrategia de vacunaci\'on afecta la din\'amica de los pesos.
\end{itemize}

Ninguna de estas tres capas exige un nuevo par\'ametro del usuario: son comportamientos emergentes de la propia estructura de la red y del estado de la epidemia.

\subsection{Glosario de m\'etricas integrado en los reportes}

Nuestros reportes PDF inclu\'ian m\'etricas (pico, $R_0$, contenci\'on, etc.) pero no explicaban qu\'e significaban ni qu\'e valores eran ``buenos''. Decidimos a\~nadir una p\'agina final de \textbf{Glosario de m\'etricas -- Gu\'ia de interpretaci\'on} (v9) con una tabla de 11 variables (12 en el reporte de lotes, agregando ``Victorias''), explicando qu\'e mide cada una y qu\'e direcci\'on es favorable desde el punto de vista de salud p\'ublica. Este glosario hace que el reporte sea autoexplicativo incluso para un lector sin formaci\'on epidemiol\'ogica.

\newpage
%==============================================================================
\section{Conclusiones y trabajo futuro}
\label{sec:conclusiones}
%==============================================================================

\subsection{Conclusiones principales}

\begin{enumerate}[leftmargin=*]
    \item \textbf{Nuestro proyecto cumple todos los puntos de la r\'ubrica del curso}: un caso de prueba principal con $N = 300$ nodos (red grande, ciudad con mezcla de estratos), carga del grafo desde archivo plano (CSV), mecanismo propio de actualizaci\'on de pesos (eventos epidemiol\'ogicos por umbral con efecto acumulativo), informaci\'on coherente con datos demogr\'aficos del DANE, y un veredicto cuantitativo claro (estrategia ganadora + score compuesto + victorias acumuladas en 60\,000 simulaciones).

    \item \textbf{Comprobamos que las estrategias estructurales superan consistentemente a la aleatoria} en la red de 300 nodos, confirmando el valor de incorporar informaci\'on topol\'ogica para decidir a qui\'enes vacunar primero.

    \item \textbf{La estrategia HUBS result\'o ganadora de forma consistente} en el experimento por lotes: con 6\,980 victorias sobre 10\,000 corridas (69{,}8\%), un pico promedio de 131 infectados y una contenci\'on del 40\%, supera a todas las dem\'as estrategias. Su fortaleza radica en bloquear los nodos de mayor alcance (hubs comunitarios de las fases 4 y 4.5) eliminando as\'i los canales de contagio inter-comunidad con un solo movimiento. La estrategia H\'ibrida (aporte propio) qued\'o cuarta (547 victorias), superada por Betweenness (1\,405) y Dijkstra Ponderado (760).

    \item \textbf{El modo por lotes (10\,000 grafos por estrategia, 60\,000 simulaciones en total) valida la robustez estad\'istica del veredicto.} Con una muestra de esta magnitud, el ranking final no es producto de una semilla favorable: HUBS 6\,980 victorias, BETWEENNESS 1\,405, Dijkstra Ponderado 760, H\'IBRIDA 547, COMUNIDADES 291, ALEATORIA 17. El margen entre el primero y el segundo (5\,575 victorias) es estad\'isticamente concluyente.

    \item \textbf{La arquitectura por capas} nos permiti\'o agregar funcionalidades sin alterar el n\'ucleo del dominio: el paso de cinco estrategias a seis (\textit{Dijkstra Ponderado}), de un modo a tres (individual, comparativo, lotes) y de exportaci\'on TXT a TXT + PDF se realiz\'o sin modificar ninguna clase del paquete \texttt{domain}.

    \item \textbf{Los eventos por umbral} (ALERTA\_LEVE, CUARENTENA, LOCKDOWN) modelan correctamente el efecto NPI de ``aplanar la curva''. En la red de 300 nodos, estos eventos se dispararon sistem\'aticamente m\'as temprano que en redes peque\~nas, confirmando la hip\'otesis H4: en redes densas, la velocidad de propagaci\'on hace que las intervenciones sean m\'as urgentes y su impacto acumulativo m\'as pronunciado.

    \item \textbf{La correci\'on metodol\'ogica} de fijar el 15\% de pacientes cero \emph{antes} de aplicar cualquier estrategia de vacunaci\'on (versi\'on 8) elimin\'o un sesgo sutil pero decisivo en la comparaci\'on. Este cambio transform\'o el sistema en un banco de pruebas estad\'isticamente v\'alido.
\end{enumerate}

\subsection{Limitaciones}

\begin{itemize}[leftmargin=*]
    \item Nuestro modelo asume \textbf{inmunidad permanente} tras la recuperaci\'on. Enfermedades reales como la influenza o el COVID-19 muestran reinfecciones que no capturamos.
    \item Calibramos los pesos $\alpha, \beta, \gamma, \delta$ de la estrategia H\'ibrida heur\'isticamente con criterio de salud p\'ublica, no mediante optimizaci\'on autom\'atica.
    \item Los factores de los eventos NPI (0.80, 0.50, 0.20) son est\'aticos y predefinidos. En la realidad, la efectividad de cada medida depende del contexto sanitario y pol\'itico. La \textbf{fatiga social} s\'i est\'a modelada (sube hasta +30\% cuando la epidemia no crece), pero la intensidad de los NPI no se ajusta din\'amicamente.
    \item No modelamos \textbf{aristas que aparecen o desaparecen} en el tiempo (nuestra red es est\'atica en estructura, solo cambian los pesos).
    \item La estrategia \textit{Betweenness} utiliza caminos m\'inimos topol\'ogicos (por saltos). Una variante que incorpore los pesos de contagio podr\'ia mejorar su efectividad a costa de mayor complejidad computacional.
\end{itemize}

\subsection{Trabajo futuro}

Como l\'ineas futuras de extensi\'on planteamos:
\begin{itemize}[leftmargin=*]
    \item \textbf{Modelo SEIRV} (con periodo de incubaci\'on Expuesto) para capturar el lapso asintom\'atico de muchas enfermedades reales.
    \item \textbf{Calibraci\'on autom\'atica de los pesos $\alpha, \beta, \gamma, \delta$} mediante algoritmos de optimizaci\'on (Random Search, Bayesian Optimization).
    \item \textbf{Betweenness ponderado}: sustituir el BFS de Brandes por Dijkstra en la Fase~1, incorporando los pesos de contagio en la definici\'on de ``camino m\'as corto''.
    \item \textbf{Aristas din\'amicas:} permitir que las relaciones aparezcan y desaparezcan en el tiempo (rotaci\'on laboral, movilidad).
    \item \textbf{Estudio de sensibilidad}: ver c\'omo cambia el ranking si la cuota de vacunaci\'on baja del 20\% al 10\% o sube al 30\%.
    \item \textbf{Comparaci\'on contra datos reales}: validar nuestras curvas $I(t)$ contra series temporales de brotes hist\'oricos (gripe AH1N1, COVID-19) ajustando $\text{probBase}$ y \texttt{diasRecuperacion}.
    \item \textbf{Calibraci\'on de la fatiga social}: actualmente los par\'ametros \texttt{UMBRAL\_FATIGA\_TURNOS} (5 turnos) y \texttt{FATIGA\_MAX} (30\%) son fijos. Una extensi\'on interesante ser\'ia ajustarlos por estrategia de vacunaci\'on o seg\'un el tipo de NPI activo.
    \item \textbf{Calibraci\'on de la vigilancia local}: el par\'ametro \texttt{VIGILANCIA\_MAX} (60\%) es uniforme para todos los nodos; en la realidad depender\'ia del estrato y la ocupaci\'on del individuo.
\end{itemize}

\newpage
%==============================================================================
\section*{Referencias}
\addcontentsline{toc}{section}{Referencias}
%==============================================================================

\begin{thebibliography}{9}

\bibitem{brandes2001}
Brandes, U. (2001).
\emph{A faster algorithm for betweenness centrality}.
Journal of Mathematical Sociology, 25(2), 163--177.

\bibitem{wattsStrogatz}
Watts, D. J., \& Strogatz, S. H. (1998).
\emph{Collective dynamics of `small-world' networks}.
Nature, 393(6684), 440--442.

\bibitem{barabasi}
Barab\'asi, A.-L., \& Albert, R. (1999).
\emph{Emergence of scaling in random networks}.
Science, 286(5439), 509--512.

\bibitem{erdosRenyi}
Erd\H{o}s, P., \& R\'enyi, A. (1960).
\emph{On the evolution of random graphs}.
Publ. Math. Inst. Hungar. Acad. Sci., 5, 17--61.

\bibitem{kermackMcKendrick}
Kermack, W. O., \& McKendrick, A. G. (1927).
\emph{A contribution to the mathematical theory of epidemics}.
Proceedings of the Royal Society A, 115(772), 700--721.

\bibitem{newmanNetworks}
Newman, M. E. J. (2010).
\emph{Networks: An Introduction}.
Oxford University Press.

\bibitem{dane}
Departamento Administrativo Nacional de Estad\'istica (DANE).
\emph{Proyecciones de poblaci\'on -- Colombia}. \url{https://www.dane.gov.co/}.

\bibitem{graphstream}
GraphStream Project.
\emph{GraphStream 2.0 -- Dynamic Graph Library}.
\url{https://graphstream-project.org/}.

\bibitem{jfreechart}
JFreeChart.
\emph{Java chart library}.
\url{https://www.jfree.org/jfreechart/}.

\bibitem{openpdf}
OpenPDF.
\emph{Java library for PDF generation (fork of iText 4)}.
\url{https://github.com/LibrePDF/OpenPDF}.

\bibitem{flatlaf}
FlatLaf.
\emph{Modern open-source cross-platform Look and Feel for Java Swing}.
\url{https://www.formdev.com/flatlaf/}.

\end{thebibliography}

\newpage
%==============================================================================
% ANEXOS
%==============================================================================
\appendix

%==============================================================================
\section{Recursos y enlaces}
\label{anexo:enlaces}
%==============================================================================

A continuaci\'on listamos los enlaces a los recursos complementarios de nuestro proyecto: el repositorio con el c\'odigo fuente completo, la grabaci\'on de la demostraci\'on en video y la presentaci\'on de diapositivas que usamos en la exposici\'on en clase.

\subsection{Repositorio de c\'odigo (GitHub)}

C\'odigo fuente completo de nuestro proyecto, archivos de datos CSV, archivos de configuraci\'on Maven y carpeta \texttt{Context/} con la documentaci\'on de dise\~no:

\begin{center}
\url{https://github.com/JoshepLizarazo/EDA_Proyecto_De_Clase_B1_2026/tree/version_3}
\end{center}

\subsection{Video de demostraci\'on (YouTube)}

Grabaci\'on en video del simulador en ejecuci\'on. Incluye la interfaz Swing, la visualizaci\'on en tiempo real con GraphStream, el disparo autom\'atico de eventos epidemiol\'ogicos, los reportes PDF que generamos y los tres modos de operaci\'on (individual, comparativo y por lotes):

\begin{center}
\url{https://youtu.be/WMGwWu-O2zY}
\end{center}

\subsection{Diapositivas de la exposici\'on}

Presentaci\'on que usamos durante la exposici\'on del proyecto en clase, con la estructura de ocho puntos que definimos como equipo:

\begin{center}
\url{https://canva.link/la8w9q5o6tz2n5o}
\end{center}

\noindent\textit{Nota:} reemplazar el placeholder por el enlace correspondiente (Canva, Google Slides, OneDrive, etc.) antes de la entrega final.

\newpage
%==============================================================================
\section{Glosario de m\'etricas (gu\'ia de interpretaci\'on)}
\label{anexo:glosario}
%==============================================================================

La tabla siguiente es una transcripci\'on del glosario que aparece como p\'agina final en los reportes PDF de nuestro simulador. Explica cada variable del modelo, qu\'e mide y qu\'e direcci\'on consideramos favorable desde el punto de vista de salud p\'ublica.

\begin{table}[H]
\centering
\small
\caption{Glosario de m\'etricas del simulador.}
\label{tab:glosario}
\begin{tabular}{l p{6cm} p{4cm}}
\toprule
\textbf{Variable} & \textbf{Qu\'e mide} & \textbf{Resultado favorable} \\
\midrule
S -- Susceptibles & Personas que pueden contagiarse & Alto al final del brote \\
I -- Infectados & Personas enfermas y contagiosas & Curva baja y estrecha \\
R -- Recuperados (afectados) & Personas que pasaron por la enfermedad & N\'umero bajo \\
V -- Vacunados & Personas inmunizadas antes de infectarse & N\'umero alto \\
Pico m\'aximo de infectados & Presi\'on sobre el sistema de salud & \textbf{Menor = mejor} \\
$t$-Pico (turno del pico) & Momento en que explot\'o el brote & \textbf{Mayor = mejor} \\
Duraci\'on del brote & Turnos hasta que $I = 0$ & \textbf{Menor = mejor} \\
Total de afectados & Cu\'antos enfermaron en total & \textbf{Menor = mejor} \\
Contenci\'on \% & \% de la poblaci\'on que NO se infect\'o & \textbf{Mayor = mejor} ($\sim$ 100\%) \\
$R_0$ estimado & Velocidad de propagaci\'on ($< 1$: se extingue) & \textbf{Menor = mejor} \\
Score compuesto $[0, 1]$ & Indicador global ponderado de las 5 m\'etricas & \textbf{Mayor = mejor} \\
Victorias (solo lote) & En cu\'antas corridas fue la mejor estrategia & \textbf{Mayor = mejor} \\
\bottomrule
\end{tabular}
\end{table}

\noindent\textit{Nota:} las m\'etricas dependen de la topolog\'ia de red espec\'ifica generada. Por eso el modo por lotes ofrece un veredicto m\'as robusto: cada estrategia se eval\'ua sobre 10\,000 grafos distintos e independientes (60\,000 simulaciones totales), lo que reduce la varianza por semilla a niveles despreciables.

\newpage
%==============================================================================
\section{Instrucciones de ejecuci\'on}
\label{anexo:instrucciones}
%==============================================================================

\subsection{Requisitos}

\begin{itemize}
    \item \textbf{Java 17} o superior (JDK con JavaFX/Swing).
    \item \textbf{Maven} 3.6 o superior.
    \item Conexi\'on a internet en la primera compilaci\'on (para descargar dependencias).
\end{itemize}

\subsection{Dependencias del proyecto}

Gestionamos las dependencias v\'ia Maven (\texttt{pom.xml}):

\begin{table}[H]
\centering
\caption{Dependencias del \texttt{pom.xml}.}
\label{tab:deps}
\begin{tabular}{l l l}
\toprule
\textbf{Librer\'ia} & \textbf{Versi\'on} & \textbf{Uso} \\
\midrule
GraphStream (\texttt{gs-core}, \texttt{gs-ui-swing}) & 2.0 & Visualizaci\'on del grafo en tiempo real \\
JFreeChart & 1.5.4 & Gr\'aficos del reporte PDF \\
OpenPDF & 1.3.34 & Generaci\'on del reporte PDF \\
FlatLaf & 3.5.4 & Tema oscuro de la UI Swing \\
\bottomrule
\end{tabular}
\end{table}

\subsection{Compilaci\'on}

Desde el directorio ra\'iz del proyecto:

\begin{lstlisting}[language=bash, basicstyle=\small\ttfamily]
mvn clean compile
\end{lstlisting}

\subsection{Ejecuci\'on}

\paragraph{Modo gr\'afico (por defecto).} Abre la interfaz Swing con FlatDarkLaf:

\begin{lstlisting}[language=bash, basicstyle=\small\ttfamily]
mvn exec:java
\end{lstlisting}

\paragraph{Modo consola.} Fuerza la interacci\'on por terminal:

\begin{lstlisting}[language=bash, basicstyle=\small\ttfamily]
mvn exec:java -Dexec.args="--consola"
\end{lstlisting}

\paragraph{Modo headless (servidor sin pantalla).} El programa detecta autom\'aticamente \texttt{GraphicsEnvironment.isHeadless()} y cae al modo consola sin intervenci\'on del usuario.

\subsection{Uso b\'asico}

\begin{enumerate}
    \item Seleccionar modo: \textbf{individual}, \textbf{comparativo}, \textbf{experimento por lotes} o \textbf{construcci\'on visual de la red}.
    \item Configurar par\'ametros: tama\~no de la red ($N = 300$ o personalizado), estrategia (solo en individual), n\'umero de grafos (solo en lotes, sin l\'imite), pausa entre fases en ms (solo en construcci\'on visual), turnos m\'aximos, d\'ias de recuperaci\'on.
    \item Opcionalmente, marcar ``Cargar red desde CSV'' para usar los archivos de \texttt{data/}.
    \item Ejecutar.
    \item Al finalizar, exportar resultados a \textbf{TXT}, \textbf{PDF} o ambos.
\end{enumerate}

\subsection{Estructura de archivos de salida}

\begin{itemize}
    \item \texttt{resultados\_simulacion.txt} -- resumen textual con la tabla comparativa y la estrategia ganadora.
    \item \texttt{reporte\_simulacion.pdf} -- reporte profesional de 8 p\'aginas (10+ p\'aginas para el modo por lotes) con portada, tabla zebra, veredicto, curvas, barras, desglose del score y glosario.
\end{itemize}

\end{document}